% concatenated from main.tex
\documentclass[12pt]{article}
\usepackage[margin=1in]{geometry}

\usepackage{graphicx} % Required for inserting images
\usepackage{amsfonts}       % blackboard math symbols
\usepackage{nicefrac}       % compact symbols for 1/2, etc.
\usepackage{microtype}      % microtypography
\usepackage{algorithm}
%\usepackage{algpseudocode}
\usepackage{algorithmic}
\usepackage{setspace}
\singlespacing

\usepackage{caption}
\usepackage{amsfonts}
\usepackage{amsmath}
\usepackage{amssymb}
\usepackage{upgreek}
\usepackage{fancyhdr}
\usepackage{titlesec}
\usepackage{indentfirst}
\usepackage{booktabs}
\usepackage{verbatim}
\usepackage{color}
\usepackage{tcolorbox}
\usepackage{mathtools}
\usepackage{bm}
\usepackage{amsthm}
\usepackage{wasysym}
\usepackage{empheq}
\usepackage{hyperref}
\usepackage{caption}
\usepackage{subcaption}
\usepackage{ulem}
\usepackage{enumerate}


\newtheorem{theorem}{Theorem}[section]
\newtheorem{corollary}{Corollary}
\newtheorem*{main}{Main Theorem}
\newtheorem{lemma}[theorem]{Lemma}
\newtheorem{proposition}[theorem]{Proposition}
\newtheorem{conjecture}{Conjecture}
\newtheorem{assumption}{Assumption}
\newtheorem*{problem}{Problem}
%\theoremstyle{definition}
\newtheorem{definition}{Definition}
\newtheorem{remark}{Remark}[section]
\newtheorem*{notation}{Notation}


%Alpha
\def\a{\alpha}
\def\k{\kappa}
\def\b{\beta}
\def\d{\delta}
\def\e{\epsilon}
\def\g{\gamma}
\def\lam{\lambda}
\def\o{\omega}
\def\p{\phi}
\def\s{\sigma}
\def\th{\theta}
\def\v{\varepsilon}
\def\th{\theta}
\def\L{\Lambda}
\def\G{\Gamma}
\def\O{\Omega}
\def\D{\Delta}
\def\p{\rho}
\def\Up{\Upsilon}
\def\Sig{\Sigma}
\def\tcr{\textcolor{black}}





%mathbb & mathcal
\def\A{{\mathbb A}}
\def\B{\mathbb B}
\def\C{\mathbb C}
\def\E{\mathbb E}
\def\N{\mathbb N}
\def\R{\mathbb R}
\def\S{{\mathbb S}}
\def\P{\mathbb P}
\def\V{\mathbb V}
\def\Z{\mathbb Z}
\def\mF{\mathcal F}
\def\mL{\mathcal{L}_{b,\Sigma}}
\def\mLhat{\mathcal{L}_{\h{b}_i,\h{\Sigma}_i}}
\def\mP{\mathcal P}
\def\mS{\mathcal S}
\def\mA{\mathcal A}
\def\mN{\mathcal N}


%left & right
\def\l{\left}
\def\r{\right}
\def\la{\left\langle}
\def\ra{\right\rangle}
\def\ll{\left\lVert}
\def\rl{\right\rVert}
\def\lv{\left\lvert}
\def\rv{\right\rvert}
\def\({\left(}
\def\){\right)}
\def\Ra{\Rightarrow}
\def\La{\Lightarrow}
\def\dsp{\displaystyle}

%cite
\usepackage[comma, sort]{natbib}
\setcitestyle{author-year}
\hypersetup{
    colorlinks=true,
    linkcolor=blue,
    filecolor=blue,      
    urlcolor=blue,
    citecolor=cyan,
}
\bibliographystyle{plainnat}

%Sign
\def\pt{\partial}
\def\nb{\nabla}
\def\cd{\cdot}
\def\ds{\displaystyle}
\def\ts{\times}
\def\qd{\quad}
\def\ol{\overline}
\def\h{\hat}
\def\t{\tilde}
\def\lles{<\!\!<}
\def\bn{\binom}

%specific
\def\dt{{\Delta t}}
\def\ds{{\Delta s}}
\def\hV{\hat{V}}
\def\tV{\t{V}}
\def\hv{\hat{v}}
\def\bV{\bar{V}}
\def\hlam{\hat{\lam}}
\def\hmu{\hat{\mu}}
\def\tmu{\t{\mu}}
\def\bmu{{\bm \mu}}
\def\hs{\hat{\Sig}}
\def\ts{\tilde{\Sig}}
\def\ts{\t{\Sig}}
\def\Lam{\Lambda}
\def\bv{{\bf v}}
\def\tcr{\textcolor{red}}
\def\tcb{\textcolor{blue}}
\def\Tau{\Upgamma}
\def\Linf{{L^\infty}}
\def\hL{\hat{L}}
\def\st{s_t}
\def\hst{\hat{s}_t} 
\def\lammin{{\lam_{\text{min}}}}
\def\lammax{{\lam_{\text{max}}}}
\newcommand{\coef}[1]{a^{(#1)}}
\def\vstar{V^*}
\def\vgstar{V^G_*}
\def\thstar{\theta_*}
\def\vg{V^G}
\def\hvi{\hat{V}_i}
\def\hvig{\hat{V}^G_i}
\def\hvigstar{\hat{V}^G_*}
\def\hvistar{\hat{V}^*_i}
\def\hE{\hat{E}}
\def\hbi{\hat{b}_i}
\def\hsigi{\hat{\Sigma}_i}
\def\rhohat{{\hat{\rho}}}
\def\drift{b}
\def\estop{\e_{\text{stop}}}
\def\thold{\th_{\text{old}}}
\def\oold{\o_{\text{old}}}
\def\Diff{\Sigma}
\def\argmax{\text{argmax}}
\title{PhiBE-Q-Learning: Bridging Off-Policy Reinforcement Learning and Continuous-Time Control}
\author{
Yuhua Zhu \thanks{Corresponding author: Department of Statistics and Data Science, University of California, Los Angeles, USA. \texttt{yuhuazhu@ucla.edu}}
\and
Yutong Ren \thanks{Department of Statistics and Data Science, University of California Los Angeles, California, USA. \texttt{ytren@ucla.edu}}
}


\date{}
\begin{document}

\maketitle
%\tcr{TO DO: 1) Prove the property and add $\max_\pi Q^\pi = Q^*$, also make the equivalent def rigorous. 2) prove coefficient difference in other three norms 3) Sample complexity results 4) Cartpole experiment}
\begin{abstract}
In this paper, we develop an off-policy method for continuous-time reinforcement learning (CTRL), where the system dynamics are governed by an unknown stochastic differential equation (SDE) and only discrete-time trajectory data are available. 
A central challenge is that the classical state--action value function $Q(s,a)$, which enables off-policy learning in discrete-time RL, does not exist in CTRL \citep{tallec2019making, jia2023q, baird1994reinforcement}. On the other hand, continuous-time control provides local notions such as the instantaneous advantage function $q(s,a)$, but these typically rely on state value function $V(s)$. To address this, we introduce a new definition of the state--action value function in CTRL and derive its governing equation.

Building on the PhiBE approximation \citep{zhu2024phibe,zhu2025optimal}, we propose iterative algorithms to approximate the optimal $Q$-function in both model-based and model-free settings using only discrete-time off-policy data. Under linear function approximation, we establish convergence guarantees and derive explicit convergence rates for the proposed method.
\end{abstract}



\section{Introduction}
%Why we should care about continoust-time RL

%What's the limitation of the current continous-time RL

%Why we follow the line of PhiBE 

%What's limitation of the current PhiBE formulation
%(why iterative and why off-policy)

%Contribution of this paper

%% related work



%Organization
Reinforcement learning (RL) has achieved remarkable success in the digital world, including AlphaGo \citep{silver2016mastering}, Atari gameplay \citep{mnih2015human}, and the fine-tuning of large language models \citep{Ziegler2019}. A key feature underlying these successes is that the system state remains static between observations. This property is natural in digital settings, making such environments well suited to RL algorithms built on the Markov decision process (MDP) framework. However, decision-making in the physical world departs from this digital paradigm in two fundamental respects. First, the objective is not merely to maximize rewards at the discrete times when data are observed, but to control performance continuously over time, including during periods without observations. In dynamic treatment regimes, for example, a patient's physiological state (e.g., glucose level) evolves continuously, while measurements are only taken intermittently. The clinical goal, however, is to regulate the state at all times \citep{emerson2023offline, cobelli2009diabetes, battelino2019clinical}. RL in the physical world thus poses a substantially more challenging objective than its classical counterpart. Second, many physical systems exhibit structured and smooth continuous-time dynamics, a feature that is largely absent in digital environments and not exploited by the MDP framework. RL in the physical world is also known as continuous-time reinforcement learning (CTRL) \citep{wang2020reinforcement}.


Off-policy learning is equally critical in real-world applications \citep{watkins1992q, sutton1998reinforcement, ma2020off}. Off-policy data consists of trajectories generated under a behavior policy that differs from the target policy, and the behavior policy may be unknown. Leveraging off-policy data can substantially improve sample efficiency, particularly in domains where data collection is costly, such as healthcare. Moreover, it enables the use of existing datasets, which is critical in settings where exploration is unsafe, such as autonomous driving.


The central question this paper studies is: what is a principled framework for CTRL when the system evolves according to an unknown stochastic differential equation (SDE), yet only discrete-time off-policy observations are available?



Existing work has largely followed two directions. The first approach seeks to reconstruct the underlying continuous-time dynamics from discrete observations and subsequently solve the induced optimal control problem using tools from stochastic control \citep{pmlr-v125-agarwal20b,pmlr-v139-yildiz21a,yang2014reinforcement,kamalapurkar2016model}. This approach has two appealing features. First, the system dynamics depend only on the state and action, therefore, it is independent of the policy, making it naturally suited for leveraging off-policy data. Second, it preserves the continuous-time structure of the problem. However, this approach also faces two fundamental limitations. First, identifying a continuous-time model from discrete observations is inherently ill-posed: in general, infinitely many continuous-time dynamics can induce the same discrete-time transition law \citep{zhu2025optimal}. As a result, model misspecification is unavoidable in practice and may lead to suboptimal policies. Second, the approach decomposes the problem into model identification and control optimization, each introducing its own source of error. When the ultimate goal is to find the optimal policy, explicitly estimating the underlying dynamics may be unnecessary and potentially inefficient.


The second approach instead reformulates CTRL as a discrete-time MDP and applies standard model-free off-policy RL algorithms, such as Q-learning or conservative Q-learning \citep{watkins1992q,kumar2020conservative}. These formulations rely solely on discrete-time transition dynamics, thereby avoiding the unidentifiability issues inherent in the first approach. In addition, model-free RL methods directly learn the optimal policy without explicitly estimating the underlying dynamics, providing a simple and broadly applicable framework across different systems. However, this approach also has two notable limitations. First, discrete-time RL algorithms do not exploit the SDE structure of the underlying dynamics. Instead, they treat all systems the same, leading to inefficient use of information and potentially suboptimal policies. More specifically, the learning error can be highly sensitive to both system parameters and reward functions \citep{zhu2025optimal}. Second, off-policy learning in the MDP framework is typically based on the action–value function $Q(s,a)$ \citep{watkins1992q}. As observed in \citep{jia2023q, tallec2019making, baird1994reinforcement}, in the continuous-time setting, the state-action $Q(s,a)$ function can be viewed as a small perturbation around the state value function $V(s)$, implying that its variation across actions is small, which makes it difficult to accurately distinguish between actions. To address this issue, these papers propose a rescaled advantage function $q(s,a)$ to amplify variation in the action space. However, this construction depends on $V$ and thus cannot be applied in the off-policy setting directly (see Section~\ref{sec:off-policy-rl} for a detailed discussion).



We tackle this challenge in two steps. First, we introduce a new definition of the state--action value function $Q$. In classical RL, a key advantage of the $Q$-function is its ability to handle off-policy data, whereas the value function $V$ does not share this property. This distinction arises because the discrete-time Bellman equation for $Q$ depends only on the transition distribution under a given action $a$, while the Bellman equation for $V$ depends on the transition induced by a policy $\pi$. Our first contribution is to define a continuous-time $Q$-function that preserves this advantage. Specifically, the proposed $Q$-function satisfies two properties: (i) it does not degenerate to the state value function $V$, unlike standard continuous-time limits of RL formulations; and (ii) its governing equation depends only on the dynamics under action $a$, not on any policy $\pi$.

The second step is to develop a new iterative algorithm to approximate the optimal value function $Q^*$. We begin with an algorithm for the setting where the underlying dynamics are known, and provide a rigorous error analysis with convergence rates under linear function approximation. Building on this, and leveraging the PhiBE framework~\citep{zhu2024phibe,zhu2025optimal}, we propose a continuous-time analogue of model-free Q-learning~\citep{watkins1992q}. This is the second key contribution of this paper. The proposed algorithm for off-policy continuous-time RL combines the strengths of continuous-time control and classical RL while avoiding their key limitations. Compared to the model-based control approaches, it relies only on discrete-time transition data and does not require explicit estimation of the underlying dynamics. In addition, it is a model-free algorithm in the sense that the update rule is invariant across different systems. Compared to standard RL algorithms, it explicitly exploits the SDE structure of the dynamics, integrating discrete-time observations with continuous-time evolution. Moreover, the resulting $Q$-function is not a vanishing perturbation of the value function $V$, which makes the formulation well-posed and numerically stable in the continuous-time limit.











%However, as noted in \citep{tallec2019making,jia2023q,  baird1994reinforcement}, in CTRL the $Q(s,a)$ function degenerates to the state value function $V(s)$ due to the infinitesimal time scale of actions. As a result, the standard mechanism that enables off-policy learning in discrete-time RL no longer applies. This degeneracy makes it difficult to exploit off-policy data in CTRL, and existing literature has not yet developed effective approaches for off-policy learning in continuous-time settings. These gaps motivate our work in this paper, where we extend PhiBE to support online updates and off-policy learning, thereby enhancing both its efficiency and practical applicability to real-world control problems.



%To address the challenge of controlling rewards continuously in time using only discrete-time data, the key is to exploit the smoothness of the underlying dynamics. \citep{zhu2024phibe,zhu2025optimal}  proposed the Physics-informed Bellman Equation (PhiBE), a framework that incorporates discrete-time information directly into the continuous-time optimal control formulation. Instead of relying on the continuous-time dynamics, PhiBE depends only on the discrete-time transition information, similar to the MDP formulation. As a result, the model is identifiable from discrete-time data and naturally leads to model-free algorithms.  Compared with standard model-free RL methods, PhiBE explicitly incorporates the smoothness structure of the underlying dynamics into the formulation, thereby reducing discretization error. These advantages makes PhiBE naturally suited for data-driven learning from trajectory samples, and it provides a principled bridge between continuous-time dynamics and discrete-time data.


%However, the algorithms proposed in the \citep{zhu2024phibe,zhu2025optimal}  are limited in two important aspects. First, hey are designed for offline learning. Second, they require on-policy data. In modern machine learning settings, it is often essential to develop online learning methods \citep{shalev2025online,zhang2023sharper}. Online learning refers to the setting where the agent improves its policy incrementally while interacting with the environment, updating its estimates after each observation or transition rather than relying on a fixed batch of data. Online learning enables agents to adapt to environments that changes over time. In addition, the model could start making good decisions earlier since updates happen incrementally, suitable for online advertisement and autonomous driving. However, existing literature has not yet explored online learning approaches tailored to continuous-time reinforcement learning (CTRL), leaving a gap in many real-world applications.


\textbf{Contributions}. We summarize the main contributions as follows.
\begin{itemize}
    \item We introduce a new state--action value function $Q$ for CTRL and derive its governing equation, which enables off-policy learning in continuous time.
    \item We develop an iterative algorithm to approximate the optimal state--action value function under known dynamics, and establish error bounds and convergence rates under linear function approximation that remain well-posed as the time discretization tends to $0$.
    \item 
    When the dynamics are unknown, we propose \textit{PhiBE-Q-Learning}, a model-free algorithm that exploits the SDE structure of CTRL problems to make effective use of off-policy data. The algorithm accommodates one-step transition data and requires parametrizing only the $Q$ function itself.
\end{itemize}

\subsection{Related work}
\subsubsection{Standard Discrete-time RL}

Classical reinforcement learning (RL) methods such as Q-learning \citep{watkins1992q}, DQN \citep{mnih2015human}, provide a powerful framework for off-policy learning in discrete-time setting. Our approach draws inspiration from these algorithms. Our update rule is an analogue of Q-learning \citep{watkins1992q} in continuous-time. %By introducing a new definition of continuous-time action-value function, we reformulate the Hamilton-Jacobi-Bellman equation into the form that only depends on the action-value function itself, disentangling it from the policy $\pi$ or state-value function $V$, thus enabling leveraging off-policy data.

In discrete-time Q-learning, convergence typically relies on the discount factor $\gamma<1$, which ensures that the optimal Bellman operator is a contraction mapping \citep{bertsekas2012dynamic}. In the continuous-time setting, the corresponding discount factor takes the form $\gamma = e^{-\beta \Delta t} \to 1$ as the time discretization $\Delta t \to 0$. This limit would, in principle, destroy the contraction property and lead to ill-conditioned or diverging guarantees. In this paper, we refine the error bounds in terms of $\Delta t$. In particular, we establish convergence and stability results that are uniform in $\Delta t$. We show that, when the underlying dynamics follow an SDE, the additional regularity of the system yields well-conditioned guarantees even $\dt\to0$. Similar results are also derived in other CTRL papers under different settings \citep{zhu2024phibe,zhu2025optimal,mou2024bellman}.

\subsubsection{Continuous-time RL(CTRL)}
A number of recent works have investigated continuous-time reinforcement learning (CTRL), where the system dynamics are modeled by SDEs \citep{jia2022policy-eval, jia2022policy-gradient, szpruch2024optimal, wang2020continuous, kim2020hamilton}.
%These studies typically design algorithms by first formulating the learning problem in continuous time and then discretizing the resulting equations to obtain implementable updates.
In particular, \citep{wang2020reinforcement} connects CTRL with stochastic control with several extensions developed in subsequent works \citep{jia2022policy-gradient, szpruch2024optimal}.
These methods typically enjoy strong theoretical guarantees and exploit the SDE structure when continuous-time data are available. However, they face three major limitations in settings with discrete-time off-policy data. First, when applied to discrete-time observations, these algorithms do not fully exploit the SDE structure of the underlying dynamics. Second, many approaches depend on long trajectories spanning from an initial state to a terminal state, which restricts applicability when only fragmented or short trajectories are available. Finally, the $q$ function defined in \citep{jia2023q,tallec2019making, baird1994reinforcement} remains dependent on the value function $V$, preventing straightforward use of off-policy data. These limitations motivate the development of methods that can learn directly from discrete observations, and leverage off-policy samples in CTRL.

%Moreover, \citep{jia2022policy-eval},  and \citep{baird1994reinforcement} pointed out that the action-value function $Q(s,a)$ cannot be directly generalized to the continuous-time setting. To address this limitation, \citep{tallec2019making} and \citep{jia2023q} introduced continuous-time analogues of the advantage function, enabling incremental value updates from sampled trajectories without requiring explicit knowledge of the system dynamics. Following this line, several algorithms based on continuous-time advantage formulations have been proposed \citep{jia2025accuracy, giegrich2024convergence, dai2023learning}.
\subsubsection{PhiBE}
To enable learning directly from discrete trajectory data, \citep{zhu2024phibe} proposed the PhiBE framework, which explicitly accounts for discretization effects and leverages the SDE structure. By doing so, PhiBE improves the accuracy and efficiency of learning from discrete trajectory data in stochastic continuous-time systems. At the same time, PhiBE still leaves certain aspects unaddressed. In particular, existing PhiBE algorithms do not fully support learning from off-policy data. This work extends the PhiBE framework to address the off-policy data utilization. Our method aims to find the solution of PhiBE instead of the original HJB equation, thereby retaining PhiBE's key advantage of handling the discrete trajectory data. Moreover, by introducing a new definition of the continuous-time action-value function $Q$, we derive an analogue of the optimal Bellman equation on $Q$ independent of policy $\pi$. This formulation allows the algorithm to utilize trajectories generated under different policies. 

Finally, we emphasize that our attention in this paper is restricted to the convergence rate when there are enough data. The finite sample error \citep{zhao2025sample,muehlebach2025sample,zhu2024phibe} is left for future work.

\subsubsection{Classical HJB solvers}


Our iterative method under known dynamics shares the same intuition as classical numerical schemes for HJB equation \citep{schaeffer2016accelerated, feng2013recent}. However, the key difference lies in the formulation and convergence mechanism. Traditional finite-difference methods approximate the HJB equation by discretizing the state space on a grid, which requires careful mesh design and often imposes CFL-type constraints on the grid spacing to ensure stability and convergence. In contrast, our approach avoids explicit spatial discretization by approximating the Galerkin projection directly from trajectory data. This eliminates the need for grid-based computation and CFL conditions, making the method more flexible and better suited for data-driven implementations in CTRL.

Moreover, based on the PhiBE framework, our method could extend to a model-free algorithm when the dynamics are unknown. Specifically, it can directly approximate the action-value function without requiring explicit knowledge of the system dynamics. Here “model-free” means that the update rule remains unchanged across different forms of underlying dynamics, allowing the algorithm to adapt automatically without model-specific modifications.

\subsection{Notation and Organization}
\paragraph{Notation} We define weighted $L^2$-norm as follows
\[
\la f, g\ra = \int f(s,a) g(s,a) \mu(s,a) dsda, \quad \ll f \rl^2 =\la f, f \ra.
\]
Here, $\mu(s,a)$ denotes the probability density of the off-policy data distribution, supported on a compact set $\Omega \subset \mathbb{R}^{d}\times \mathbb{R}^{m}$.

For vector $x\in \R^n$, and function $f(s,a)$, we define $\ll \cdot \rl_2, \ll \cdot \rl_\infty$ as follows,
\[
\ll x \rl^2_2 = \sum_{i=1}^n x_i^2, \quad \ll f \rl_\infty = \sup_{(s,a)\in\O} \ll f \rl_2
\]

If no specific remarks are given, then $\nb, \nb^2$ refer to $\nb_s, \nb^2_s$. 



\paragraph{Organization.}
Section~\ref{sec:preliminary} introduces the problem setup. 
Section~\ref{sec:off-policy-rl} presents the new state--action $Q$-function. 
Section~\ref{sec:off-policy-known} develops the algorithm under known dynamics and establishes its theoretical guarantees, while Section~\ref{sec:off-policy-phibe} derives the model-free algorithm for unknown dynamics. 
Section~\ref{sec:4.3} reports numerical experiments that validate the theory.



\section{Preliminary}\label{sec:preliminary}
\subsection{Stochastic optimal control}
We consider a continuous-time stochastic optimal control problem defined on a filtered probability space $(\Omega, \mathcal{F}, \{\mathcal{F}_t\}, \mathbb{P})$. The state process $\{s_t\}_{t \ge 0}$ takes values in $\mathbb{R}^d$ and evolves according to a controlled stochastic differential equation (SDE):
\begin{equation}\label{dyn}
    ds_t = b(s_t, a_t)\, dt + \sigma(s_t, a_t)\, dB_t,
\end{equation}
where $b: \mathbb{R}^d \times \mathcal{A} \to \mathbb{R}^d$ is the drift term, $\sigma: \mathbb{R}^d \times \mathcal{A} \to \mathbb{R}^{d \times k}$ is the diffusion term, and $B_t$ is a $k$-dimensional standard Brownian motion.  
At time $t$, the agent takes an action $a_t =\pi(s_t)$ according to a feedback control. Here the feedback control is defined as a deterministic mapping 
$\pi : \mathbb{R}^d \to \mathcal{A}$ from the state space to a compact action space $\mathcal{A} \subseteq \mathbb{R}^m$.  
Given a fixed policy $\pi$ and initial state $s_0 = s$, the value function is defined as
\begin{equation}\label{def of V-pi}
    V^\pi(s) = \mathbb{E}\left[ \int_0^\infty e^{-\beta t} r(s_t, \pi(s_t))\, dt \;\middle|\; s_0 = s \right],    
\end{equation}
where $\beta > 0$ is the discount rate. The goal of stochastic optimal control problem is to find the optimal feedback control that maximizes the value function,
\begin{equation}\label{sop-goal}
    \pi^*(s) = \argmax_\pi\, V^\pi(s), \quad V^*(s) = V^{\pi^*}(s).
\end{equation}

We assume the following assumption to ensure the well-posedness of the above stochastic control problem. 
\begin{assumption}\label{ass}
The action space $\mA$ is compact; the functions $b(s,a)$, $\sigma(s,a)$, and $r(s,a)$ are continuous in $a$, and locally uniformly Lipschitz continuous in $s$. The reward function $r(s,a)$ is uniformly bounded.
\end{assumption}
Moreover, the optimal value function $V^*(s)$ and the optimal feedback policy $\pi^*$ defined in \eqref{sop-goal}  satisfy
\[
\beta V^*(s) = \sup_{a\in\mA}[r(s,a) + (\mL V^*)(s,a)],\quad 
\pi^*(s) = \arg\sup_{a\in\mA}[r(s,a) + (\mL V^*)(s,a)],
\]
Here
\begin{equation}
\mL = b(s,a)\cdot \nabla_s + \frac{1}{2}\Sigma(s,a) :\nabla_s^2\label{eq:PE}
\end{equation}
with $\Sigma(s,a) = \sigma(s,a)\sigma^\top(s,a) \in \R^{d\times d}$ and $\Sigma(s,a) : \nabla_s^2 = \Sigma_{i,j}\Sigma(s,a)_{ij}\partial_{s_i}\partial_{s_j}$. We refer readers to \citep{FS06, YZ99, P09} for detailed accounts of the classical stochastic control theory.




\subsection{Reinforcement learning}
In the setting of reinforcement learning, the dynamics of the environment, i.e., $b(s,a), \sigma(s,a)$ are unknown.
We assume access to discretized information, either in the form of transition densities
\[
    \rho_{\Delta t}(s' \mid s,a), \qquad \mathbb{R}^d \times \mathbb{R}^d \times \mathcal{A} \to [0,\infty),
\]
which denotes the probability density function of $s_{\Delta t}$ given $s_0 = s$ and $a(\tau) = a$ for $\tau \in (0,\Delta t)$, or a dataset of transitions
\begin{equation}\label{data}
\{(s_i, a_i, s'_i, r_i)\}_{i=1}^n, \qquad (s_i, a_i)\sim \mu(s,a), \quad s'_i \sim \rho_{\Delta t}(\,\cdot \mid s_i,a_i), \quad r_i = r(s_i, a_i).    
\end{equation}
Here the actions $\{a_i\}_{i=1}^n$ are generated by some (possibly unknown) behavior policy and are not controlled by the learner. Note that we consider a piecewise constant control where the action is frozen at the left endpoint, i.e., $a_t = a_0$ for $t \in [0,\Delta t]$ in the data set, while the goal is still to find the optimal feedback control $\pi^*(s)$ defined in \eqref{sop-goal}. 


\subsection{PhiBE}\label{sec:2.2}
As indicated in \citep{zhu2024phibe,zhu2025optimal}, the optimal-PhiBE provides a better approximation to the continuous-time optimal control problem compared to the Bellman equation when the underlying dynamics are smooth or the reward function oscillates, as it preserves the underlying SDE structure. Therefore, in this work, we adopt PhiBE to approximate the value function. The first-order Optimal-PhiBE corresponding to \citep{zhu2025optimal} is given by 
\begin{definition}[PhiBE]
\begin{equation}\label{optimal-phibe}
  \beta \hat{V}^*(s)
=
\sup_{a \in \mathcal{A}}
\left\{
    r(s,a)+
     \mathcal{L}_{\hat{b}, \hat{\Sigma}} \hat{V}^*(s)
\right\},  
\end{equation}
where
\begin{equation}\label{phibe-bsigma}
\begin{aligned}
\hat{b}(s,a)&=
\frac{1}{\Delta t}
\int (s' - s)\rho_\dt(s'|s,a)ds'\\
\hat{\Sigma}(s,a)
&=
\frac{1}{\Delta t}
\int(s'- s)(s' - s)^\top \rho_\dt(s'|s,a)ds'
\end{aligned}    
\end{equation}
\end{definition}
As shown in Theorem 3.5 of \citep{zhu2025optimal}, under suitable regularity conditions on the dynamics $b$ and $\sigma$, the optimal policy $\hat{\pi}^*(s)$ derived from PhiBE achieves an $O(\dt)$ approximation to the true optimal policy, when evaluated in terms of the true value function under the respective policies. We will use the PhiBE formulation in Section \ref{sec:off-policy-phibe} when only discrete-time information are available.


\section{PhiBE Q-Learning}\label{sec:off-policy-rl}

In this section, we aim to find a method that leverages the off-policy data. Off-policy data refers to experience collected under a behavior policy that is different from the target policy being optimized, in which the behavior policy may be unknown. Off-policy learning plays a crucial role in continuous-time reinforcement learning (CTRL). By allowing the use of data generated under arbitrary behavior policies, it significantly improves sample efficiency and makes it possible to reuse trajectories, which is especially important when collecting new data is costly. Furthermore, in many practical applications such as healthcare or autonomous driving, exploration is inherently risky. Off-policy learning provides a way to leverage existing datasets without requiring unsafe interactions. It also helps reduce the limitation caused by using data only from the current policy, enabling learning to reach an optimal policy. Finally, the ability to separate safe data collection from optimal policy optimization highlights the practical necessity of off-policy methods for CTRL, motivating our extension of PhiBE to this setting.

In standard RL, the key reason why $Q(s,a)$ naturally supports off-policy learning is because the Bellman equation 
\begin{equation}\label{discrete-time be}
 Q^\pi(s,a) = r(s,a) + \gamma \, \E_{s'\sim \rho(\cdot|s,a)} \l( \E_{a'\sim\pi(\cdot|s) }\l[ Q^\pi(s',a') \r]\r)   
\end{equation}
is determined by the unknown transition dynamics $\rho(\cdot|s,a)$ independent of the target policy $\pi$. While the Bellman equation for the value functions $V^\pi(s)$
\[
V^\pi(s) = r^\pi(s) + \gamma \, \E_{s'\sim \rho^\pi(\cdot|s)} \l( V^\pi(s') \r) 
\]  
is determined by the unknown marginal transition kernel  $\rho^\pi(s'|s) \;=\; \int \pi(a|s)\,\rho(s'|s,a)\, da$, which explicitly depends on $\pi$. As a result, when data are generated from a different behavior policy, the estimation of $V^\pi$ is inherently biased if the off-policy data is directly applied. In contrast, the Bellman optimality equation for $Q^*(s,a)$ is conditioned directly on $(s,a)$, and the transition kernel $\rho(s'|s,a)$ depends only on the environment dynamics, not on $\pi$. Thus, any sample $(s,a,s',r)$ collected under an arbitrary behavior policy provides an unbiased estimate of the right-hand side of the Bellman equation, making off-policy learning feasible for $Q$.  

Although the Bellman equation for $V^\pi(s)$ can, in principle, be adapted to off-policy learning by applying importance sampling to correct for the mismatch between the target policy $\pi$ and the behavior policy, this approach faces two challenges in practice. First, importance sampling often leads to high variance, especially in long-horizon problems or when the behavior and target policies differ substantially. Second, the behavior policy itself may be unknown, making it difficult to compute the required importance sampling ratios. These limitations motivate the use of $Q$-functions, whose Bellman equation is naturally off-policy without requiring such corrections.

A direct generalization of the discrete-time action-value function $Q^\pi(s,a)$ to the continuous-time setting is not feasible. As proved in \citep{jia2023q,tallec2019making}, the discrete-time action-value function $Q_{\Delta t}^\pi(s,a)$ defined as
\begin{equation}\label{def:q-dt}
Q_{\Delta t}^\pi(s,a) = 
\E_{\substack{
a_i \sim \pi(\cdot|s_i), i\geq1\\
s_{i+1} \sim \rho_{\Delta t}(\cdot \mid s_i, a_i), i\geq 0
}} \Big[\sum_{i = 0}^\infty\gamma^ir(s_i,a_i)|s_0 = s, a_0 = a \Big]    
\end{equation}
will converge to $V^\pi$ as $\Delta t \to 0$. The reason is that, as the discretization step $\Delta t \to 0$, the duration over which an action influences the environment also vanishes. In this limit, the effect of choosing action $a$ at state $s$ becomes negligible, and the action-value function $Q^\pi(s,a)$ effectively degenerates to the value function $V^\pi(s)$. 

In \citep{jia2023q, tallec2019making}, the authors propose an action-value function defined as
\begin{align*}
    q^\pi(s,a) := r(s,a) + [\mL V^\pi](s,a) - \beta V^\pi(s).
\end{align*}
This formulation, however, depends explicitly on $V^\pi$, whose estimation typically relies on on-policy transitions generated by $\pi$. We note that both \citet{jia2023q} and \citet{tallec2019making} propose off-policy algorithms for jointly estimating $q^\pi$ and $V^\pi$. Nevertheless, these approaches have two limitations when applied to discrete-time data. First, they require parameterizing two separate functions, $q^\pi$ and $V^\pi$, thereby introducing two sources of approximation error. Second, the resulting algorithms with discrete-time data do not fully exploit the smooth SDE structure. 

Therefore, in order to handle the off-policy data within one value function and embed the SDE structure, we introduce a modified definition of the action-value function in the CTRL framework.





In Sections~\ref{sec:4.1.1} and \ref{sec:off-policy-known}, we begin by analyzing the known-dynamics setting, where the continuous transition dynamics $(b, \Sigma)$ are assumed to be known. Moreover, we derive an iterative algorithm and establish the exponential convergence of the proposed scheme. Section~\ref{sec:4.2.1} further demonstrates that PhiBE allows us to extend the analysis from known dynamics to the unknown-dynamics case, where only discrete transition dynamics $\rho_{\Delta t}(s'|s,a)$ are available. Finally, Section~\ref{sec:4.2.2} also presents a model-free, off-policy data-driven algorithm that implements the proposed method when the dynamics are unknown and only discrete-time data is available, followed by numerical experiments in Section~\ref{sec:4.3} to demonstrate the performance of the approach.


\subsection{Continuous-time state-action value function $Q(s,a)$}\label{sec:4.1.1}
\begin{definition}[Continuous-time Bellman equation for $Q^\pi(s,a)$]\label{def of Q-pi}
\begin{equation}\label{pde of q}
Q^\pi(s,a) = r(s,a) + \mL \l[\E_{a\sim\pi} (Q^\pi)\r](s,a) + (1-\beta) \E_{a\sim\pi} [Q^\pi(s,a)]
\end{equation}

\iffalse
When policy $\pi(s) \in C^2$ is deterministic, then an alternative way for continuous-time $Q^\pi(s,a)$ is 
\begin{equation}\label{pde of q-deter-pi}
\beta Q^\pi(s,a) = r(s,a) + b(s,a)\cdot D_\pi Q(s,a) + \frac12\Sigma(s,a):H_\pi Q (s,a)
\end{equation}
where 
\[
\begin{aligned}
    D_\pi Q(s,a) = &M_\pi(s)^\top\nb_{(s,a)} Q = \nabla_s Q(s,a)+\nabla_s\pi(s)^\top \nabla_a Q(s,a) \\
H_\pi Q(s,a) = &M_\pi(s)^\top\nabla_{s,a}^2Q(s,a)M_\pi(s)
+
\sum_{k=1}^m
\partial_{a_k}Q(s,a)\nabla_s^2\pi_k(s)
\end{aligned}
\]
and $M_\pi(s)
=[
I_d,
\nabla_s\pi(s)]^\top.$

\fi
\end{definition}
Several remarks on the definition of $Q^\pi$ in \eqref{pde of q} are in order.

First, comparing \eqref{pde of q} with the continuous-time Bellman equation for $V^\pi$ defined in \eqref{def of V-pi},
\[
\beta V^\pi(s) = r^\pi(s) + \mathcal{L}_{b^\pi,\Sigma^\pi} V^\pi(s), 
\quad b^\pi(s) = \E_{a\sim\pi}[b(s,a)], 
\quad \Sigma^\pi(s) = \E_{a\sim\pi}[\Sigma(s,a)],
\]
we see a key structural difference. The operator in \eqref{pde of q} is governed by the dynamics $(b(s,a), \Sigma(s,a))$, and therefore does not depend on the target policy $\pi$. This mirrors the discrete-time Bellman equation for $Q_{\Delta t}$ in \eqref{discrete-time be}. In contrast, the dynamics for $V^\pi$ depend on $(b^\pi, \Sigma^\pi)$ and thus on $\pi$. Consequently, when the dynamics are unknown, $Q^\pi$ naturally accommodates off-policy data, whereas $V^\pi$ does not.

Second, the continuous-time $Q^\pi$ can also be viewed as the sum of the value function and the instantaneous advantage function:
\[
Q^\pi(s,a) = \lim_{\Delta t\to 0} \left[\frac{Q^\pi_{\Delta t}(s,a) - V^\pi_{\Delta t}(s)}{\Delta t}\right] + V^\pi(s),
\]
where $Q^\pi_{\Delta t}$ is defined in \eqref{def:q-dt} and $V^\pi_{\Delta t}(s) = \E_{a\sim\pi}[Q^\pi_{\Delta t}(s,a)]$. As shown in \citep{jia2023q}, instantaneous advantage function is given by
\begin{equation}\label{def of q}
q^\pi(s,a) := r(s,a) + [\mathcal{L} V^\pi](s,a) - \beta V^\pi(s).
\end{equation}
Accordingly, $Q^\pi$ can be written as
\begin{equation}\label{eq-Q-2}
Q^\pi(s,a) = r(s,a) + [\mathcal{L} V^\pi](s,a) - \beta V^\pi(s) + V^\pi(s).
\end{equation}
Taking expectation over $a \sim \pi$ yields
\[
\E_{a\sim\pi}[Q^\pi(s,a)] 
= r^\pi(s) + \mathcal{L}_{b^\pi,\Sigma^\pi} V^\pi(s) - \beta V^\pi(s) + V^\pi(s)
= V^\pi(s),
\]
where the last equality follows from the Bellman equation for $V^\pi$. Substituting $V^\pi(s) = \E_{a\sim\pi}[Q^\pi(s,a)]$ into \eqref{eq-Q-2} recovers the continuous-time Q-Bellman equation \eqref{pde of q}.


Third, we adopt the formulation in Definition \ref{def of Q-pi}, rather than \eqref{eq-Q-2}, because it expresses the equation entirely in terms of the $Q$-function. This avoids the need to explicitly compute $V^\pi$ and is essential for  off-policy learning.

\iffalse
Several remarks regarding the second definition of $Q^\pi$ in \eqref{pde of q-deter-pi} are in order. First, if one inserts $a = \pi(s)$ in \eqref{pde of q-deter-pi}, one has,
\[
\beta Q^\pi(s,\pi(s)) = r(s,\pi(s) ) + b(s,\pi(s))\cdot \nb_s Q(s,\pi(s)) + \frac12 \Sigma(s,\pi(s)) :\nb_s^2Q^\pi(s,\pi(s))
\]
Let $\hat{V}^\pi(s) =  Q^\pi(s,\pi(s))$, then $\hat{V}(s)$ satisfies the same equation as $V^\pi(s)$ defined in \eqref{def of V-pi}. As proved in 
\fi





Next, we define the continuous-time optimal Bellman equation for $Q^*(s,a)$ as follows.
\begin{definition}[Continuous-time optimal Bellman equation for $Q^*(s,a)$]\label{def of Q-pde}
\begin{equation}\label{pde of Qstar}
    \beta Q^*(s,a) = r(s,a) + \mL \big[\max_{a'} Q^*(s,a')\big] .
\end{equation}
\end{definition}

\iffalse
To verify that \eqref{pde of Qstar} is consistent with the standard RL relation $V^*(s) = \max_{a} Q^*(s,a)$, we provide a formal proof below.
Define the operators
\begin{align}\label{operator-Q}
    T_Q(Q) &= r(s,a) + \mL \big[\max_{a} Q(s,a)\big] + (1-\beta)\,\big[\max_{a} Q(s,a)\big],\\
T_V(V) &= \max_{a} \Big(r(s,a) + [\mL V](s,a) + (1-\beta)V\Big).
\end{align}
Let $Q^*$ denote the fixed point of $T_Q$, i.e., $T_Q(Q^*) = Q^*$. Define $\hat{V}(s) = \max_{a} Q^*(s,a)$. Then
\begin{align*}
    T_V(\hat{V}) 
    &= \max_{a}\Big(r(s,a) + [\mL \hat{V}](s,a) + (1-\beta)\hat{V}\Big) \\
    &= \max_{a}\Big(r(s,a) + \mL \big[\max_{a} Q^*(s,a)\big] + (1-\beta)\,\big[\max_{a} Q^*(s,a)\big]\Big) \\
    &= \max_{a} Q^*(s,a) = \hat{V}.
\end{align*}
Hence, $\hat{V}$ is a fixed point of the operator $T_V$. Since $V^*$, the optimal value function, is also the unique fixed point of $T_V$ (by the existence and uniqueness of the solution to the HJB equation), we conclude that $\hat{V} = V^*$. Consequently, the identity $V^*(s) = \max_{a} Q^*(s,a)$ holds.
\fi

To verify that \eqref{pde of Qstar} is consistent with the standard RL relation $V^*(s) = \max_{a} Q^*(s,a)$, we provide a formal proof below.
Define the operators
\begin{align}\label{operator-Q}
T_Q(Q) &=\frac{1}{\beta}\l( r(s,a) + \mL \big[\max_{a} Q(s,a)\big] \big]\r),\\
T_V(V) &= \frac1\beta\max_{a} \Big(r(s,a) + [\mL V](s,a) \Big).
\end{align}
Let $Q^*$ denote the fixed point of $T_Q$, i.e., $T_Q(Q^*) = Q^*$. Define $\hat{V}(s) = \max_{a} Q^*(s,a)$. Then
\begin{align*}
    T_V(\hat{V}) 
    &= \frac{1}{\beta}\max_{a}\Big(r(s,a) + [\mL \hat{V}](s,a)\Big) = \max_{a}\frac1\beta\Big(r(s,a) + \mL \big[\max_{a'} Q^*(s,a')\big]\Big) \\
    &= \max_{a} Q^*(s,a) = \hat{V}.
\end{align*}
Hence, $\hat{V}$ is a fixed point of the operator $T_V$. Since $V^*$, the optimal value function, is also the unique fixed point of $T_V$ (by the existence and uniqueness of the solution to the HJB equation), we conclude that $\hat{V} = V^*$. Consequently, the identity $V^*(s) = \max_{a} Q^*(s,a)$ holds.







We summarize several key properties of the continuous-time $Q^\pi$ and $Q^*$ as follows:
\begin{itemize}

\item The function $Q^\pi(s,a)$ admits the equivalent definition
\[
Q^\pi(s,a) = \lim_{\Delta t \to 0} \left[\frac{Q^\pi_{\Delta t}(s,a) - V^\pi_{\Delta t}(s)}{\Delta t}\right] + V^\pi(s).
\]

\item $Q^\pi(s,a)$ induces the same policy improvement as $q^\pi(s,a)$, under both the greedy and softmax policy update rules,
\[
\arg\max_a q^\pi(s,a) = \argmax_a Q^\pi(s,a), \quad \pi(a|s) = \frac{\exp( \lambda q^\pi(s,a))}{Z_q(s)} = \frac{\exp( \lambda Q^\pi(s,a)) }{Z_Q(s)}
\]
where $Z_q, Z_Q$ are the normalizing constant.
\item The solution to the continuous-time $Q$-Bellman equation \eqref{pde of q} is consistent with the value function $V^\pi$ in \eqref{def of V-pi}, in the sense that
\[
\E_{a \sim \pi}[Q^\pi(s,a)] = V^\pi(s).
\]

\item The solution to the continuous-time optimal $Q$-Bellman equation \eqref{pde of Qstar} satisfies
\[
\max_{a} Q^*(s,a) = V^*(s) = \max_{\pi} V^\pi(s) = \max_{\pi} \E_{a \sim \pi}[Q^\pi(s,a)].
\]

\end{itemize}



\subsection{Q-learning under Known Dynamics}\label{sec:off-policy-known}
In this section, we introduce an iterative algorithm to solve  \eqref{pde of Qstar} for $Q^*(s,a)$ under linear bases. Let $\Phi(s): \mathcal{S} \to \mathbb{R}^n$ be a set of basis functions. We approximate the solution by $\hat{Q}(s,a) = \Phi(s,a)^\top \th^*$ using Galerkin method. The Galerkin approximation $\hat{Q}(s,a)$ satisfies, 
\begin{equation}\label{galerkin-solu}
    \la r(s,a) + \mL \l(\max_a \l[\Phi(s,a)^\top\th^*\r]\r) + (1-\beta) \max_a\l[\Phi(s,a)^\top\th^*\r] - \Phi(s,a)^\top\th^*,\Phi(s,a)\ra= 0
\end{equation}
\begin{equation}\label{galerkin-solu}
    \la r(s,a) + \mL \l(\max_a \l[\Phi(s,a)^\top\th^*\r]\r) -\beta \Phi(s,a)^\top\th^*,\Phi(s,a)\ra= 0
\end{equation}

However, directly solving the above equation for $\th^*$ is infeasible. We add a time dependence on $\theta(t)$, and derive an ODE for $\th(t)$
\begin{equation*}
\begin{aligned}
 \frac{d}{dt}\theta(t) 
=&
\big\langle 
r(s,a) 
+ \mL \Big(\max_a \big[\Phi(s,a)^\top \theta(t)\big]\Big) - \beta\Phi(s,a)^\top \theta(t),
\ \Phi(s,a)
\big\rangle       
\end{aligned}
\end{equation*}
Discretize the time evolution, one has
\begin{equation}\label{eq:newalgoite}
\begin{aligned}
\theta_{n+1} =& \theta_n  + \alpha_n \la  r(s,a) + \mL \Big(\max_a [\Phi(s,a)^\top \theta_n]\Big)-\beta\Phi(s,a)^\top \theta_n, \ \Phi(s,a) \ra
\end{aligned}
\end{equation}

We establish the properties of the proposed algorithm in two steps. First, Theorem~\ref{thm:galerkin} shows that, under Assumption~\ref{ass-conv}, the Galerkin solution $\theta^*$ defined in \eqref{galerkin-solu} provides a good approximation to the true value function $Q^*(s,a)$ defined in \eqref{pde of Qstar}. Second, Theorem~\ref{thm:galerkin_convergence} shows that the iterates $\theta_n$ generated by Algorithm~\eqref{eq:newalgoite} converge to the Galerkin approximation $\theta^*$.

\begin{assumption}\label{ass-conv}
\begin{itemize}
\iffalse
    \item [(A1)] The smallest eigenvalue $\lam_1$ of the Gram matrix $G = \int \Phi\Phi^\top \mu dsda$  is positive. The largest eigenvalue $\lam_2$ of the Gram matrix is bounded, and the bases $\ll\Phi\rl_\infty$ are bounded. In addition, the largest eigenvalue of Gram matrix $G_\nb, G_{\nb^2}$ of $\nb\Phi$, $\nb^2\Phi$ are bounded. Let $c_3 := \max_{v\in \mathbf{S}^{n-1}}\frac{\|\max_a |v^\top \Phi|\|}{\|\Phi^\top v\|}$, we assume that 
    \begin{equation}\label{def of cbeta}
        c_\beta = 1-c_3|1-\beta| > 0
    \end{equation}
    \item[(A2)] The functions $ b(s,a)\in C^1(\mathbb{S} \times \mathbb{A}),\Sigma(s,a) \in C^2(\mathbb{S} \times \mathbb{A})$ satisfy: 
    \begin{align*}
        \|\nabla b\|_{\infty} + \frac{1}{2}\|\nabla^2 \Sigma\|_{\infty} &\leq \frac{c_\beta}{6c_3}, \quad 
        \|b\|_{\infty} + \|\nabla \Sigma\|_{\infty} \leq \frac{c_\beta c_1}{6c_3}, \quad 
         \frac{1}{2}\|\Sigma\|_{\infty} \leq \frac{c_\beta c_2}{6c_3}.
    \end{align*}
        where $c_1 := \min_{v\in \mathbf{S}^{n-1}}\frac{\|\Phi^\top v\|}{\|\nabla\Phi^\top v  
           \|}, c_2 := \min_{v\in \mathbf{S}^{n-1}}\frac{\|\Phi^\top v\|}{\|\nabla^2\Phi^\top v  
           \|}$
\fi 
     \item [(A1)] The smallest eigenvalue $\lam_1$ of the Gram matrix $G = \int \Phi\Phi^\top \mu dsda$  is positive. The largest eigenvalue $\lam_2$ of the Gram matrix is bounded, and the bases $\ll\Phi\rl_\infty$ are bounded. In addition, the largest eigenvalue of Gram matrix $G_\nb, G_{\nb^2}$ of $\nb\Phi$, $\nb^2\Phi$ are bounded.
     \item[(A2)] $\| b\|_{C^1} , \|\Sigma\|_{C^2} \leq L$
     \item[(A3)]    Let $c_0 := \max_{v\in \mathbf{S}^{n-1}}\frac{\|\max_a |v^\top \Phi|\|}{\|\Phi^\top v\|}$, $c_1 := \max_{v\in \mathbf{S}^{n-1}}\frac{\|\nabla\Phi^\top v  \|}{\|\Phi^\top v\|}, c_2 := \max_{v\in \mathbf{S}^{n-1}}\frac{\|\nabla^2\Phi^\top v  
           \|}{\|\Phi^\top v\|}$. We assume that 
    \begin{equation}\label{def of cbeta}
        \beta > C_{L,\Phi}: = c_0 L (1+c_1+c_2)
    \end{equation}

\end{itemize} 
\end{assumption}

Several remarks on the assumptions are in order. Assumption (A1) constrains the choice of basis functions $\Phi(s)$, Assumption (A2) imposes conditions on the underlying dynamics, and Assumption (A3) requires the discount parameter $\beta$ to be large enough.

In Assumption (A1), the requirement that the smallest eigenvalue of the Gram matrix be strictly positive is equivalent to linear independence of the basis functions in the weighted space $L^2(\mu)$. 


In Assumption (A3), we note that the constants $c_0, c_1, c_2$ are always finite under Assumption (A1). For $c_0$, since $\ll v^\top \Phi \rl = v^\top G v \geq \lam_1$ for $\ll v \rl_2 = 1$, one has $c_0 \le \frac{\|\Phi\|_\infty}{\sqrt{\lambda_1}}.$ For $c_1, c_2$,  let $\lambda_3$ and $\lambda_4$ denote the largest eigenvalues of the matrices $G_{\nabla}$ and $G_{\nabla^2}$, respectively. Then one can further bound $c_1 \leq \frac{\lambda_3}{\lambda_1}, \; c_2 \leq \frac{\lambda_4}{\lambda_1}$.


The discounted coefficient $\beta$ and the constant $C_{L,\Phi}$ are the key quantities that governing the performance of the algorithm. Later in Theorems \ref{thm:galerkin} and \ref{thm:galerkin_convergence},  we will show that both the approximation accuracy and convergence behavior improve as $\beta$ becomes larger or $C_{L,\Phi}$ becomes smaller. Since $C_{L,\Phi}$ depends on $c_0, c_1, c_2$
through $L$, and smaller $c_0, c_1, c_2$ weakens the requirement on $\beta$
while improving algorithmic performance, we discuss in the following the
conditions under which $c_0, c_1, c_2$ are small.


\paragraph{Uniform exploration of the behavior policy controls $c_0$.} Let $\mu(s) = \int_\A \mu(s,a) da$ be the marginal distribution over states, and define the conditional
distribution $\pi(a|s)$ s.t. $\mu(s,a) = \mu(s)\pi(a|s)$, which corresponds to the behavior policy generating the data. We claim that the uniform behavior policy $\pi^u(a|s) = 1/|\mathcal{A}|$ minimizes the constant $c_0$ in the worst case over all
bounded bases $\Phi$.
Any other policy $\pi \neq \pi^u$ leaves some region of the action space
systematically under-covered relative to $\pi^u$, and a basis functions can
always be constructed that places its informative directions precisely in
that region, inflating $c_0(\pi)$ above $c_0(\pi^u)$.

The precise statement is presented in the following lemma, and the proof is given in Appendix \ref{sec:proof of uniform-optimal}.
\begin{lemma}\label{lemma:uniform-optimal}
The uniform policy $\pi^u$ is the unique minimax optimizer to
\begin{equation*}
\pi^u = \operatorname*{argmin}_{\pi \in \Pi}\;
\max_{\Phi \in \mathcal{F}_B}\; c_0(\pi, \Phi)^2.
\end{equation*}
where $\Pi := \Bigl\{\pi(\cdot|\cdot) : \pi(a|s) \geq 0,\;
\int_\mathcal{A}\pi(a|s)\,da = 1 \quad \mu(s)\text{-a.e.}\Bigr\}$, $\mathcal{F}_B := \{\Phi \text{ continuous} : \|\Phi\|_\infty \leq B\}$.    
\end{lemma}

The minimax formulation in Lemma~\ref{lemma:uniform-optimal} is necessary
for the following reason. For a fixed known basis functions $\Phi$, a policy concentrating on actions
where $\Phi$ has high energy can achieve a smaller $c_0$ than $\pi^u$.
However, such a policy requires prior knowledge of the geometry of $\Phi$, precisely the object the offline RL algorithm is trying to learn.
The minimax formulation $\min_{\pi} \max_{\Phi} c_0(\pi, \Phi)^2$ captures
the natural requirement that the behavior policy must be chosen without this
knowledge. Under this criterion, $\pi^u$ is uniquely optimal: it is the only
policy that cannot be adversarially exploited by any choice of basis functions,
because it covers every action equally regardless of where $\Phi$ places its
informative directions.


\paragraph{Smoothness and orthonormality of the bases $\Phi$ controls $c_1$ and $c_2$.}
The constants $c_1$ and $c_2$ are relative smoothness constants of the basis
functions $\Phi$. They are controlled by two independent
mechanisms. First, the numerators $\|\nabla\Phi^\top v\|$ and
$\|\nabla^2\Phi^\top v\|$ are small when $\Phi$ has small derivatives across
$\mathcal{S}\times\mathcal{A}$,  which can be achieved by regularizing the derivatives, for instance through smoothness penalties or bounded weight norms
in a neural network parameterization. Second, the denominator
$\min_v\|\Phi^\top v\| = \lam_1$ is large when the
basis functions are orthonormal under $\mu$, so that no projection direction is starved of signal. Any
deviation from orthogonality under $\mu$ reduces $\lam_1$
and increases $c_1$ and $c_2$.


\iffalse
Assumption (A2) imposes smoothness and boundedness conditions on the drift $b$, diffusion $\Sigma$. Specifically, the bounds require that the magnitude and gradient magnitude of $b$ and $\Sigma$ are sufficiently small compared with $\lam_1$, so that the Galerkin approximation remains stable. In addition, we note that the constants $c_1, c_2, c_3$ are always strictly positive under Assumption (A1). This is because $\ll v^\top \Phi \rl = v^\top G v \geq \lam_1$ for $\ll v \rl_2 = 1$. Let $\lambda_3$ and $\lambda_4$ denote the largest eigenvalues of the matrices $G_{\nabla}$ and $G_{\nabla^2}$, respectively. Then one can further bound $c_1 \geq \frac{\lambda_1}{\lambda_3}>0, \; c_2 \geq \frac{\lambda_1}{\lambda_4}>0$. These bounds highlight that the admissible magnitudes of $\|b\|_{C^1}$ and $\|\Sigma\|_{C^2}$ depend on the sensitivity of the basis functions. When $\Phi$ is smooth and dominated by low-frequency components under $L^2(\mu)$, larger dynamics can be accommodated. In contrast, rapidly varying bases (with large gradients or curvature) require correspondingly smaller drift and diffusion to maintain stability.









In Assumption (A1), the requirement that the smallest eigenvalue of the Gram matrix be strictly positive is equivalent to linear independence of the basis functions in the weighted space $L^2(\mu)$. The condition $c_\beta>0$ requires either a sufficiently small $c_3$ or the discount coefficient $\beta$ sufficiently close to one.




The constant $c_\beta$ is one of the key quantities governing both the assumptions on the dynamics and the performance of the algorithm. When $c_\beta$ is larger, Assumption (A2) imposes weaker restrictions on the dynamics, allowing for more oscillatory systems. Later, in Theorems \ref{thm:galerkin} and \ref{thm:galerkin_convergence}, we will show that both the approximation accuracy and convergence behavior improve when $c_\beta$ is larger. Additional discussion will be provided after Theorem \ref{thm:galerkin_convergence}. 

Next, we discuss the regime $\beta\in(0,1)$. One observes a positive relationship between the discount coefficient $\beta$ and the assumptions imposed on the dynamics. When $\beta$ is larger, i.e., closer to $1$, Assumption (A2) permits more oscillatory systems. In contrast, when $\beta$ is closer to $0$, the optimal control problem becomes less discounted, and Assumption (A2) requires stronger regularity conditions on the dynamics. This type of trade-off is also common in infinite-horizon discounted optimal control problems. However, our assumption requires $\beta$ to remain within a moderate range and not be too far from $1$. This assumption is mainly technical and may potentially be relaxed with a more refined analysis, which we leave for future work.






We next discuss the constant $c_3$. First, note that $c_3$ is always finite. Indeed, for any $v\in\mathbb R^d$ with $\|v\|_2=1$, $\|v^\top \Phi\|^2 = v^\top G v \ge \lambda_1,$ where $\lambda_1$ is the smallest eigenvalue of $G$. Hence, $c_3 \le \frac{\|\Phi\|_\infty}{\sqrt{\lambda_1}}.$

Next, consider the special case where the action space $\mathcal A$ is finite with $K=|\mathcal A|$ actions. Let
$\mu(s)=\sum_{i=1}^K \mu(s,a_i)$ be the marginal distribution over states, and define the conditional distribution $\mu(a_i\mid s)=\frac{\mu(s,a_i)}{\mu(s)},$
which corresponds to the behavior policy generating the data. Define $\mu_0(s)=\min_i \mu(a_i\mid s).$ Then,
\begin{equation}\label{ineq-3}
\l(\max_{a_i} | v^\top \Phi(s,a_i) |\r)^2 \leq \sum_i | v^\top \Phi(s,a_i) |^2 \leq \frac{1}{\mu_0(s)}\sum_i | v^\top \Phi(s,a_i) | \mu(a_i|s)    
\end{equation}
Therefore,
\[
\begin{aligned}
&\|\max_a |v^\top \Phi|\|^2 = \int \l(\max_{a_i} | v^\top \Phi |\r)^2 \mu(s) ds \\
\leq& \frac{1}{\mu_0(s)} \int \sum_i | v^\top \Phi(s,a_i) | \mu(a_i|s)\mu(s) ds = \max_s\l(\frac{1}{\mu_0(s)}\r) \ll v^\top \Phi \rl^2   , 
\end{aligned}
\]
which implies 
\begin{equation}\label{ineq-2}
    c_3\leq \sqrt{\max_s\l(\frac{1}{\mu_0(s)}\r)}.
\end{equation} 
Note that the RHS of \eqref{ineq-2} is minimized when the behavior policy is uniform over the action space, i.e., $\mu(a_i\mid s)=1/K$, in which case $c_3 \leq \sqrt{K}$. %More generally, suppose $\mu(a_i\mid s)\ge c_b/K$ for some constant $c_b>0$. Then $c_3\le \sqrt{K/c_b}$, which implies that a more imbalanced behavior policy leads to a larger $c_3$ and hence a smaller $c_\beta$. 
This suggests that when the behavior policy explores the action space more uniformly, the proposed algorithm achieves faster convergence and better approximation accuracy in the presence of model error. Finally, we note that the bound $\sqrt{K}$ is generally loose because the first inequality in \eqref{ineq-3} can be highly conservative. In practice, the actual value of $c_3$ is often substantially smaller than $\sqrt{K}$ in this special case.








Assumption (A2) imposes smoothness and boundedness conditions on the drift $b$, diffusion $\Sigma$. Specifically, the bounds require that the magnitude and gradient magnitude of $b$ and $\Sigma$ are sufficiently small compared with $\lam_1$, so that the Galerkin approximation remains stable. In addition, we note that the constants $c_1, c_2, c_3$ are always strictly positive under Assumption (A1). This is because $\ll v^\top \Phi \rl = v^\top G v \geq \lam_1$ for $\ll v \rl_2 = 1$. Let $\lambda_3$ and $\lambda_4$ denote the largest eigenvalues of the matrices $G_{\nabla}$ and $G_{\nabla^2}$, respectively. Then one can further bound $c_1 \geq \frac{\lambda_1}{\lambda_3}>0, \; c_2 \geq \frac{\lambda_1}{\lambda_4}>0$. These bounds highlight that the admissible magnitudes of $\|b\|_{C^1}$ and $\|\Sigma\|_{C^2}$ depend on the sensitivity of the basis functions. When $\Phi$ is smooth and dominated by low-frequency components under $L^2(\mu)$, larger dynamics can be accommodated. In contrast, rapidly varying bases (with large gradients or curvature) require correspondingly smaller drift and diffusion to maintain stability.

It is worth emphasizing that our assumptions are primarily for theoretical guarantees and are not sharp. 
In the numerical experiments, we show that even when these conditions are violated, the algorithm still exhibits stable convergence in practice. 

\fi


Before we present the theorem, we first prove a Lemma that will be frequently used in the theoretical guarantees. 


\iffalse
\begin{lemma}\label{lemma:ineq}
Under Assumption \ref{ass-conv}, for any $Q = \Phi^\top \th$ in the linear space spanned by $\Phi$, one has the following upper bound,
\[
\ll (1-\beta)Q +\mL^*Q \rl \leq  \frac{1-c_\beta/2}{c_3}\ll Q \rl, \quad \text{where}\quad \mL^* f = \nb\cdot\l[ bf + \frac12\nb\cdot (\Sigma f)\r] .
\]
\end{lemma}
\begin{proof}
First note that
\[
\begin{aligned}
&\ll (1-\beta)Q +\mL^*Q \rl\\
\leq & |1-\beta| \ll Q \rl + \l(\ll \nb b\rl_\infty + \frac12\ll\nb^2 \Sigma\rl_\infty \r) \ll Q\rl + \l( \ll b \rl_\infty + \ll \nb\Sigma\rl\r)  \ll \nb Q\rl + \frac12\ll\Sigma\rl_\infty \ll \nb^2 Q  \rl \\
\leq& \frac{1-c_\beta}{c_3} \ll Q \rl + \frac{c_\beta}{6c_3}\l(\ll Q  \rl + c_1\ll \nb\Phi^\top \th\rl + c_2 \ll \nb^2\Phi^\top \th\rl \r)\\
\leq& \frac{1-c_\beta}{c_3} \ll Q \rl + \frac{c_\beta}{2c_3}\ll Q  \rl  = \frac{1-c_\beta/2}{c_3} \ll Q \rl
\end{aligned}    
\]
where the first inequality is due to Cauchy–Schwarz inequality; the second inequality is by Assumption (A2) and definition of $c$; the third inequality is because $Q = \th^\top\Phi$ and for $\forall 0\neq \th\in \R^n$, one has
\[
\begin{aligned}
&c_1\ll \nb\Phi^\top \th \rl = \min_{v\in \mathbf{S}^{n-1}}\frac{\|\Phi^\top v\|}{\|\nabla\Phi^\top v \|} \ll  \nb\Phi^\top \th \rl \leq \ll \th^\top \Phi \rl, \quad\\
&c_2\ll  \nb^2\Phi^\top \th \rl = \min_{v\in \mathbf{S}^{n-1}}\frac{\|\Phi^\top v\|}{\|\nabla^2\Phi^\top v\|}\ll  \nb^2\Phi^\top \th \rl \leq \ll \th^\top \Phi \rl .    
\end{aligned}
\] 
\end{proof}
\fi 




\begin{lemma}\label{lemma:ineq}
Under Assumption \ref{ass-conv}, for any $Q = \Phi^\top \th$ in the linear space spanned by $\Phi$, one has the following upper bound,
\[
\ll \mL^*Q \rl \lesssim  L(1+c_1 + c_2)\ll Q \rl, \quad \text{where}\quad \mL^* f = \nb\cdot\l[ bf + \frac12\nb\cdot (\Sigma f)\r] .
\]
where the inequality holds up to a uniform constant that is independent of the problem. 
\end{lemma}
\begin{proof}
First note that
\[
\begin{aligned}
\ll \mL^*Q \rl
\leq &  \l(\ll \nb b\rl_\infty + \frac12\ll\nb^2 \Sigma\rl_\infty \r) \ll Q\rl + \l( \ll b \rl_\infty + \ll \nb\Sigma\rl\r)  \ll \nb Q\rl + \frac12\ll\Sigma\rl_\infty \ll \nb^2 Q  \rl \\
\lesssim& L \l(\ll Q  \rl + \ll \nb\Phi^\top \th\rl + \ll \nb^2\Phi^\top \th\rl \r) 
\leq L(1+c_1+c_2)\ll Q \rl
\end{aligned}    
\]
where the first inequality is due to Cauchy–Schwarz inequality; the second inequaltiy is by Assumption (A2); the third inequality is because $Q = \th^\top\Phi$ and for $\forall 0\neq \th\in \R^n$, one has
\[
\begin{aligned}
&\ll \nb\Phi^\top \th \rl =  \frac{\ll \nb\Phi^\top \th \rl }{\ll \Phi^\top \th \rl}  \ll \Phi^\top \th \rl \leq \max_{v\in \mathbf{S}^{n-1}} \frac{\ll \nb\Phi^\top v \rl }{\ll \Phi^\top v \rl}  \ll \Phi^\top \th \rl = c_1 \ll \th^\top \Phi \rl,
\end{aligned}
\] 
and similar inequality can be obtained for $\ll \nb^2\Phi^\top \th\rl \leq c_2 \ll Q \rl$.
\end{proof}


\color{black}
\begin{theorem}[Galerkin solution accuracy]\label{thm:galerkin}
Under Assumption \ref{ass-conv}, the Galerkin solution $\hat{Q}^*(s,a) = \Phi(s,a)^\top \th^*$ that satisfies \eqref{galerkin-solu} has an accuracy of
\[
\ll \hat{Q}^* - Q^* \rl \leq \frac{\beta + C_{L,\Phi}}{\beta - C_{L,\Phi}}  \min_{\th}\ll Q^* -\Phi^\top\th \rl
\]
where $C_{L,\Phi} = Lc_0(1+c_1+c_2)$ is a constant only depending on $\Phi, b, \Sigma$.
\end{theorem}

Several remarks on the above theorem are in order. First, when the optimal value function $Q^*(s,a)$ can be represented by the bases $\Phi(s,a)$, then the Galerkin solution is accurate. Second, one can view $\min_{\th} \ll Q^* - \Phi^\top \th \rl$ as the model error, that is the best approximation one can obtain in the space spanned by the bases $\Phi$. The theorem shows that the Galerkin solution $\hat{Q}^*$ achieves, up to a multiplicative constant, the smallest possible error among all linear combinations of $\Phi$. Therefore, whenever the linear basis $\Phi(s,a)$ provides a good approximation to $Q^*$, the Galerkin solution inherits this accuracy and is guaranteed to be near-optimal within the chosen function class. 


In the presence of model approximation error, the accuracy of the Galerkin
solution improves as $\beta$ becomes larger or $C_{L,\Phi}$ becomes smaller.
Based on the discussion following Assumption~\ref{ass-conv}, this implies
that the Galerkin approximation becomes more accurate under the following
conditions: 1) $\beta$ is larger; 2) the behavior policy explores the action space more uniformly; 3) the basis functions $\Phi$ are smoother; 4) the basis functions are more orthogonal under the off-policy data distribution.
%The Galerkin method does not introduce a significant additional error beyond the inherent approximation error of the basis. 

\begin{proof}
We divide the error $e(s,a) = Q^*(s,a) - \hat{Q}^*(s,a)$ into two parts, 
\begin{equation}\label{error decomp}
    e(s,a) = e_G(s,a) + e_P(s,a), \quad e_P = Q^* - \Phi^\top u, \quad e_G = \Phi^\top v, \quad v = u - \th^*
\end{equation}
where $u, v\in \R^d, u = \arg\min_{\th}\ll Q^* -\Phi^\top\th \rl$. Here $e_G$ represents the Galerkin error, and $e_P$ represents the model error.  






Taking inner product of the continuous-time Q-function \eqref{pde of Qstar} with $e_G$, one has
\[
\la r + \mL(\max_a Q^*), e_G\ra = \la \beta Q^*, e_G\ra
\]
Taking inner product of $v$ with the Galerkin equation \eqref{galerkin-solu}, one has
\[
\la r + \mL(\max_a \hat{Q}^*) , e_G\ra = \la  \beta \hat{Q}^*, e_G\ra
\]
subtracting the two equations gives,
\[
\la \mL(\max_a Q^* - \max_a \hat{Q}^*) , e_G\ra = \beta \ll e_G \rl^2 + \beta\la e_P, e_G\ra
\]
Now the LHS can be bounded by
\[
\begin{aligned}
LHS =& \la \max_a Q^* - \max_a \hat{Q}^*,  \mL^*e_G\ra \leq L(1+c_1+c_2)\ll \max_a Q^* - \max_a \hat{Q}^* \rl  \ll e_G \rl   \\
\leq &  L(1+c_1+c_2) \ll \max_a |e_G  + e_P| \rl \ll e_G \rl \\
\leq &  L(1+c_1+c_2)\l(\ll \max_a |e_G| \rl\ll e_G \rl       +\ll \max_a|e_P| \rl \ll e_G \rl \r)
\leq Lc_0(1+c_1+c_2)\l(\ll e_G \rl^2       + \ll e_P \rl \ll e_G \rl \r)
\end{aligned}
\]
where the first equality is due to integration by parts, and the first inequality is by Lemma \ref{lemma:ineq}. The last inequality is because of Assumption \ref{ass-conv}/(A1), for any $e(s,a) =e^\top \Phi(s,a) $ with $\ll e \rl_2\neq 0$, 
\begin{equation}\label{ineq-1}
    \ll \max_a|e| \rl = \frac{\ll \max_a|e| \rl}{\ll e \rl} \ll e \rl \leq \max_{v\in \mathbf{S}^{n-1}}\frac{\|\max_a |v^\top \Phi|\|}{\|\Phi^\top v\|} \ll e \rl  = c_0\ll  e \rl.
\end{equation}
Therefore, denote $C_{L,\Phi} = Lc_0(1+c_1+c_2)$, one has
\[
(\beta - C_{L,\Phi})\ll e_G \rl^2 \leq  C_{L,\Phi}\ll e_P \rl \ll e_G \rl -\beta  \la e_P, e_G\ra \leq(\beta +C_{L,\Phi})\ll e_P \rl \ll e_G \rl
\]
which implies 
\[
\ll e_G \rl \leq \frac{\beta + C_{L,\Phi}}{\beta - C_{L,\Phi}}\ll e_P \rl.
\]
One completes the proof by inserting the definition of $e_P$ in \eqref{error decomp}. 



\end{proof}

The convergence guarantee is derived in the following theorem.

\iffalse
\begin{theorem}[Convergence of Q-learning under known dynamics]
\label{thm:galerkin_convergence}
Under Assumption \ref{ass-conv}, the Q-learning iteration defined in \eqref{eq:newalgoite} with learning rate $\alpha_n = \alpha \in (0, \frac{c_\beta\lam_1}{4\lam_2^2})$ will converges to the Galerkin solution $ \theta^* $ defined in \eqref{galerkin-solu} in the following rate,
\[
\ll \th_n  - \th^* \rl_2 \leq (1-\alpha c_\beta \lam_1 + 4\alpha^2\lam_2^2)\ll \th_0 - \th^* \rl_2
\]
By setting the learning rate $\alpha = \frac{c_\beta\lam_1}{8\lam_2^2}$, the error decays as below:
\[
\|\theta_n - \theta^*\|_2 \leq \l(\sqrt{1 - \l(\frac{c_\beta\lam_1}{4\lam_2}\r)^2}\r)^n \|\theta_0 - \theta^*\|_2.
\]
\end{theorem}
\fi


\begin{theorem}[Convergence of Q-learning under known dynamics]
\label{thm:galerkin_convergence}
Under Assumption \ref{ass-conv}, the Q-learning iteration defined in
\eqref{eq:newalgoite} with learning rate $\alpha_n = \alpha \in (0, \frac{2m}{K^2})$
converges to the Galerkin solution $\theta^*$ defined in \eqref{galerkin-solu}
at a linear rate. By setting the learning rate
$\alpha = \frac{m}{K^2}$, it achieves the optimal rate,
\[
\|\theta_n - \theta^*\|_2 \leq \left(\sqrt{1 - \frac{m^2}{K^2}}\right)^n \|\theta_0 - \theta^*\|_2,
\]
where $m = \lambda_1(\beta - C_{L,\Phi}), \quad K = \lambda_2(\beta + C_{L,\Phi})$.
\end{theorem}


%Several remarks on the theorem are in order.  The optimal convergence rate depends explicitly on the spectrum of the Gram matrix through $\lambda_1$ and $\lambda_2$, as well as the key parameter $\beta, C_{L,\Phi}$. In particular, better-conditioned bases lead to faster convergence. In the special case where the basis functions are orthonormal under $L^2(\mu)$, the Gram matrix becomes the identity, so $\lambda_1 = \lambda_2 = 1$, yielding the optimal rate. Therefore the convergence rate becomes faster under the same condition that makes the Galerkin solution more accurate, that is 1) $\beta$ is larger; 2) the behavior policy explores the action space more uniformly; 3) the basis functions $\Phi$ are smoother; 4) the basis functions are more orthogonal under the off-policy data distribution.

Several remarks on the theorem are in order. The optimal convergence rate
depends explicitly on the spectrum of the Gram matrix through $\lambda_1$
and $\lambda_2$, as well as the key parameters $\beta$ and $C_{L,\Phi}$.
In particular, better-conditioned bases lead to faster convergence: in the
special case where the basis functions are orthonormal under $L^2(\mu)$,
the Gram matrix becomes the identity, so $\lambda_1 = \lambda_2 = 1$,
yielding the fastest possible rate. The convergence rate therefore becomes
faster under the same conditions that make the Galerkin solution more
accurate, namely: 1) $\beta$ is larger; 2) the behavior policy explores the
action space more uniformly; 3) the basis functions $\Phi$ are smoother;
4) the basis functions are more orthogonal under the off-policy data
distribution.





\def\mH{\mathcal{H}}

\iffalse
\begin{proof}
Define
\begin{equation}
        \mathcal{H}(\theta) = \left\langle \Phi^T \theta - (\mathcal{L}_{b,\Sigma} + (1-\beta)) \max_a \Phi^T \theta - r, \Phi \right\rangle.
\end{equation}   
\paragraph{Monotonicity.}
We first show that $\mH$ is strongly monotone. 
Let $Q = \th^\top\Phi, P = \eta^\top \Phi$, expanding the left-hand side, we have
\begin{align*}
&\l(\mathcal{H}(\theta) -  \mathcal{H}(\eta) \r)^\top (\theta - \eta)\\
= &\|Q-P\|^2 - (1-\beta)\langle \max_a Q - \max_a P, Q-P \rangle - \la (b \nabla + \frac{1}{2} \Sigma: \nabla^2)(\max_a Q - \max_a P), Q-P \ra\\
=&\|Q-P\|^2 - \langle \max_a Q - \max_a P, ( (1-\beta) (Q-P) +\mL^*(Q-P) \ \rangle \\
\geq& \|Q-P\|^2 -\frac{1-c_\beta/2}{c_3}\ll\max_a Q - \max_a P\rl \ll Q - P \rl  \geq \frac{c_\beta}2\|Q-P\|^2
\end{align*}
where the first inequality is by Lemma \ref{lemma:ineq}, and the second inequality is obtained by applying \eqref{ineq-1}, 
\[
\ll\max_a Q - \max_a P\rl \leq \ll \max_a |Q - P| \rl \leq c_3\ll Q - P \rl
\]
Again by Assumption (A1), one has
\[
\ll Q - P \rl^2 = (\th - \eta)^\top\l[\int\Phi\Phi^\top \mu(s,a)dsda\r](\th - \eta) \geq \lam_1\ll \th - \eta \rl_2^2 
\]
Therefore, one has
\[
\l(\mathcal{H}(\theta) -  \mathcal{H}(\eta) \r)^\top (\theta - \eta) \geq \frac{c_\beta\lam_1}{2}\ll \th - \eta\rl_2^2.
\]

\paragraph{Lipschitz Continuity of $ \mathcal{H} $.}
We verify that $ \mathcal{H} $ is Lipschitz. Observe
\begin{align*}
&\|\mathcal{H}(\theta) - \mathcal{H}(\eta)\| 
\leq \ll \langle \Phi^\top (\theta-\eta), \Phi \rangle\rl + \ll  \max_aQ - \max_a P \rl \ll (1-\beta)\Phi + \mL^*\Phi \rl\\
\leq & G(\th - \eta) + \frac{1-c_\beta}{2} \ll Q - P \rl\ll \Phi\rl 
\leq 2\lam_2 \ll \th - \eta \rl_2 
\end{align*}
where Lemma \ref{lemma:ineq} and \eqref{ineq-1} is applied to obtained the second inequality, and Assumption \ref{ass-conv}/(A1) is used in the third inequality. Then one has, 
\begin{align*}
&\|\mathcal{H}(\theta) - \mathcal{H}(\eta)\| 
\leq 2\lam_2\ll \th - \eta \rl_2
\end{align*}

\paragraph{Contraction Property.}
For fixed learning rate $\alpha_k = \alpha$, since the Galerkin iteration update is
\[
\theta_{k+1} = \theta_k - \alpha_k \mathcal{H}(\theta_k),
\]
and $\mathcal{H}(\th^*) = 0$, one has
\[
\begin{aligned}
&\ll \th_{n+1} - \th^* \rl_2^2 = \ll \th_n -\th^* - \alpha (\mathcal{H}(\th_n) -  \mathcal{H}(\th^*)) \rl_2^2 \\
\leq&     \|\theta_n - \th^*\|_2^2 - 2\alpha  (\mathcal{H}(\theta_n) - \mathcal{H}(\th^*))^\top (\theta_n - \th^*)+ \alpha^2 \|\mathcal{H}(\theta_n) - \mathcal{H}(\th^*)\|_2^2
\end{aligned}
\]
By Monotonicity and Lipschitz continuity of $\mathcal H$, we have
\begin{align*}
\ll \th_{n+1} - \th^* \rl_2^2\leq (1 - 2m\alpha + \alpha^2 K^2) \|\theta_n - \th^*\|_2^2,\quad m = \frac{\lam_1c_\beta}{2}, \quad K = 2\lam_2.
\end{align*}
First note that since $K>m$, one has $1 - 2\alpha m+ \alpha^2 K^2 > 0$. Second, one chooses $0< \alpha < \frac{2m}{K^2}$ then the mapping becomes a contraction. Lastly, by setting $\alpha = m/K^2$, the convergence rate is optimized as 
\[
\ll \th_n - \th^* \rl_2 \leq \l(\sqrt{1-\frac{m^2}{K^2}}\r)^n \ll\th_0 - \th^* \rl_2 
\]. 
\end{proof}
\fi




\begin{proof}
Define
\begin{equation}
        \mathcal{H}(\theta) = \left\langle \beta \Phi^T \theta - \mathcal{L}_{b,\Sigma}\l(  \max_a \Phi^T \theta\r) - r, \Phi \right\rangle.
\end{equation}   
\paragraph{Monotonicity.}
We first show that $\mH$ is strongly monotone. 
Let $Q = \th^\top\Phi, P = \eta^\top \Phi$, expanding the left-hand side, we have
\begin{align*}
&\l(\mathcal{H}(\theta) -  \mathcal{H}(\eta) \r)^\top (\theta - \eta)\\
= &\beta\|Q-P\|^2 - \la (b \nabla + \frac{1}{2} \Sigma: \nabla^2)(\max_a Q - \max_a P), Q-P \ra\\
=&\beta\|Q-P\|^2 - \langle \max_a Q - \max_a P, \mL^*(Q-P) \ \rangle \\
\geq&\beta \|Q-P\|^2 - L(1+c_1+c_2)\ll\max_a Q - \max_a P\rl \ll Q - P \rl  \geq (\beta - C_{L,\Phi})\|Q-P\|^2
\end{align*}
where the first inequality is by Lemma \ref{lemma:ineq}, and the second inequality is obtained by applying \eqref{ineq-1}, 
\[
\ll\max_a Q - \max_a P\rl \leq \ll \max_a |Q - P| \rl \leq c_0\ll Q - P \rl
\]
Again by Assumption (A1), one has
\[
\ll Q - P \rl^2 = (\th - \eta)^\top\l[\int\Phi\Phi^\top \mu(s,a)dsda\r](\th - \eta) \geq \lam_1\ll \th - \eta \rl_2^2 
\]
Therefore, one has
\[
\l(\mathcal{H}(\theta) -  \mathcal{H}(\eta) \r)^\top (\theta - \eta) \geq \lam_1(\beta - C_{L,\Phi})\ll \th - \eta\rl_2^2.
\]

\paragraph{Lipschitz Continuity of $ \mathcal{H} $.}
We verify that $ \mathcal{H} $ is Lipschitz. Observe
\begin{align*}
&\|\mathcal{H}(\theta) - \mathcal{H}(\eta)\| 
\leq \beta\ll \langle \Phi^\top (\theta-\eta), \Phi \rangle\rl + \ll  \max_aQ - \max_a P \rl \ll  \mL^*\Phi \rl\\
\leq & \beta \ll G(\th - \eta)  \rl + C_{L,\Phi}\ll Q - P \rl\ll \Phi\rl 
\leq(\beta+ C_{L,\Phi}) \lam_2  \ll \th - \eta \rl_2 
\end{align*}
where Lemma \ref{lemma:ineq} and \eqref{ineq-1} is applied to obtained the second inequality, and Assumption \ref{ass-conv}/(A1) is used in the third inequality. 

\paragraph{Contraction Property.}
For fixed learning rate $\alpha_k = \alpha$, since the Galerkin iteration update is
\[
\theta_{k+1} = \theta_k - \alpha_k \mathcal{H}(\theta_k),
\]
and $\mathcal{H}(\th^*) = 0$, one has
\[
\begin{aligned}
&\ll \th_{n+1} - \th^* \rl_2^2 = \ll \th_n -\th^* - \alpha (\mathcal{H}(\th_n) -  \mathcal{H}(\th^*)) \rl_2^2 \\
\leq&     \|\theta_n - \th^*\|_2^2 - 2\alpha  (\mathcal{H}(\theta_n) - \mathcal{H}(\th^*))^\top (\theta_n - \th^*)+ \alpha^2 \|\mathcal{H}(\theta_n) - \mathcal{H}(\th^*)\|_2^2
\end{aligned}
\]
By Monotonicity and Lipschitz continuity of $\mathcal H$, we have
\begin{align*}
\ll \th_{n+1} - \th^* \rl_2^2\leq (1 - 2m\alpha + \alpha^2 K^2) \|\theta_n - \th^*\|_2^2,\quad m = \lam_1(\beta - C_{L,\Phi}), \quad K = \lam_2(\beta + C_{L,\Phi}).
\end{align*}
First note that since $K>m$, one has $1 - 2\alpha m+ \alpha^2 K^2 > 0$. Second, one chooses $0< \alpha < \frac{2m}{K^2}$ then the mapping becomes a contraction. Lastly, by setting $\alpha = m/K^2$, the convergence rate is optimized as 
\[
\ll \th_n - \th^* \rl_2 \leq \l(\sqrt{1-\frac{m^2}{K^2}}\r)^n \ll\th_0 - \th^* \rl_2 
\]. 
\end{proof}




\subsection{Q-learning under unknown dynamics}\label{sec:off-policy-phibe}
In the previous section, we assumed that the continuous dynamics $b, \sigma$ are known. In this section, we assume $b, \sigma$ are unknown. In Section \ref{sec:4.2.1}, we introduce PhiBE to naturally apply the algorithm and theoretical guarantees when only the discrete-time transition dynamics are known. In \ref{sec:4.2.2}, we describe how to implement the update rule \eqref{eq:newalgoite} in a data-driven manner when the transition dynamics is unknown and only the state-action-next-state-reward tuples  $\{(s_i, a_i,s'_i,r_i)\}_{i=1}^n$ as defined in \eqref{data} are available.

\subsubsection{PhiBE setting}\label{sec:4.2.1}
When only the discrete-time transition dynamics $\rho_{\Delta t}(s'|s, a)$ are available, one can replace the unknown drift and diffusion $b,\Sigma$ in the equation \eqref{pde of Qstar} for $Q^*$ with the approximated $\hat{b}$ and $\hat{\Sigma}$ according to \eqref{phibe-bsigma}, and one ends up with the Optimal-PhiBE-Q equation

\iffalse
\begin{definition}[Optimal-PhiBE-Q]\label{def:phibe-q}
    \begin{equation}\label{pde of Qhatstar}
    \hat{Q}^*(s,a) = r(s,a) + \mathcal{L}_{\hat{b},\hat{\Sigma}} \big[\max_{a} \hat{Q}^*(s,a)\big] + (1-\beta)\,\big[\max_{a} \hat{Q}^*(s,a)\big], 
\end{equation}
where $\hat{b}, \hat{\Sigma}$ are defined in \eqref{phibe-bsigma}. 
where $\hat{b}, \hat{\Sigma}$ are defined in \eqref{phibe-bsigma}. 
\end{definition}
\fi

\begin{definition}[Optimal-PhiBE-Q]\label{def:phibe-q}
\begin{equation}\label{pde of Qhatstar}
    \beta \hat{Q}^*(s,a) = r(s,a) + \mathcal{L}_{\hat{b},\hat{\Sigma}} \big[\max_{a} \hat{Q}^*(s,a)\big] , 
\end{equation}
where $\hat{b}, \hat{\Sigma}$ are defined in \eqref{phibe-bsigma}. 
\end{definition}

Substituting $(\hat{b}, \hat{\Sigma})$ into \eqref{operator-Q} yields a corresponding operator, whose fixed point corresponds to the optimal PhiBE solution to \eqref{optimal-phibe}.
Solving the Optimal-PhiBE-Q equation in Definition \ref{def:phibe-q}, the same algorithm \eqref{eq:newalgoite} applies with the approximated $\hat{b}, \hat{\Sigma}$. Therefore as long as $\hat{b}, \hat{\Sigma}$ satisfies the Assumption (A2), the corresponding algorithm 
\begin{equation}\label{eq:newalgoite-phibe}
\begin{aligned}
\theta_{n+1} =& \theta_n  + \alpha_n \bigg\langle  r(s,a) + \mathcal{L}_{\hat{b},\hat{\Sigma}}  \Big(\max_a [\Phi(s,a)^\top \theta_n]\Big)-\beta \Phi(s,a)^\top \theta_n, \ \Phi(s,a) \bigg\rangle 
\end{aligned}
\end{equation}
recovers the same limiting behavior as established in Theorem \ref{thm:galerkin} and Theorem \ref{thm:galerkin_convergence}.

\begin{corollary}
Under Assumption \ref{ass-conv}/(A1), (A3). For $\forall \epsilon>0$, if $\ll b \rl_{C^1}, \ll \Sigma \rl_{C^2} \leq L-\epsilon$, then there exists $\dt$ small enough, s.t., the Galerkin solution $\t{Q}^*(s,a) = \Phi(s,a)^\top\t{\th}^*$ for Optimal-PhiBE-Q,
\begin{equation}\label{galerkin-solu-phibe}
    \la r(s,a) + \mathcal{L}_{\hat{b},\hat{\Sigma}}  \l(\max_a \l[\Phi(s,a)^\top \t{\th}^*\r]\r)  - \beta\Phi(s,a)^\top\t{\th}^*,\Phi(s,a)\ra= 0
\end{equation}
has an accuracy of 
\[
\ll \t{Q}^*(s,a) - \hat{Q}^*(s,a) \rl \leq \frac{\beta + C_{L,\Phi}}{\beta - C_{L,\Phi}}\min_\th \ll \hat{Q}^*(s,a) - \Phi^\top \th \rl_\infty.
\]
In addition, with the same learning rate as in Theorem \ref{thm:galerkin_convergence}, the iterative algorithm \eqref{eq:newalgoite-phibe} will converges to $\t{\th}^*$ in the same rate as in Theorem \ref{thm:galerkin_convergence}.
\end{corollary}
\begin{proof}
As proved in \citep{zhu2025optimal}, $\ll \nb^k(\hat{b} - b) \rl_\infty \leq C\dt$ with $k = 0,1$, $\ll \nb^k(\Sigma - \hat{\Sigma})\rl_\infty \leq C\dt$ with $k = 0,1,2$. Therefore, for sufficiently small $\dt$, one has 
\[
\|b - \hat{b}\|_{C^1} +\| \Sigma - \hat{\Sigma}\|_{C^2} \leq \e
\]
then by the assumption, one has
$\|\hat{b}\|_{C^1} \leq \|b\|_{C^1} + \|b - \hat{b}\|_{C^1}  \leq L.$ Similarly, one can prove with the same $\dt$, $\|\hat{\Sigma}\|_{C^2} \leq \|\Sigma\|_{C^2} + \| \Sigma - \hat{\Sigma}\|_{C^2}  \leq L.$

Therefore, the approximated $\hat{b}, \hat{\Sigma}$ satisfies Assumption \ref{ass-conv}/(A2), then the guarantees in Theorems \ref{thm:galerkin} and \ref{thm:galerkin_convergence} follows. 

\end{proof}

\subsubsection{Data-driven algorithm}\label{sec:4.2.2}
When only discrete-time off-policy data \eqref{data} is available, our goal is to approximate the optimal action-value function $Q^*(s,a)$. 
%The algorithm is presented in Algorithm~\ref{algo2}.


Following the algorithm \eqref{eq:newalgoite-phibe} introduced in last section, we approximate the optimal-Q function $Q^*(s,a)$ using the bases $\Phi(s,a)=(\phi_1,\dots,\phi_n)^\top$.  In each iteration, given $Q_n(s,a) = \Phi(s,a)^\top\th_n$, one first computes 
\[
V_n(s) = \max_a Q_n(s,a).
\]
Then a mini-batch $\mathcal{B}_n\subset\mathcal{B}$ is sampled, and for each tuple $(s_i,a_i,r_i)$ in the mini-batch, the TD-error is computed as
\begin{equation}\label{td-data}
 \delta(s_i,s_i',a_i,r_i) = r_i + \mathcal{L}_{b_i, \Sigma_i}V_n(s_i)  - \beta Q_n(s_i,a_i),   
\end{equation}
which serves as an unbiased estimate of the PhiBE residual. Here $\mathcal{L}_{b,\Sigma}$ is defined in \eqref{eq:PE}, and 
\[
b_i = \frac{s'_i - s_i}{\Delta t}, \quad 
\hat \Sigma_i = \frac{(s'_i - s_i)(s'_i - s_i)^\top}{\Delta t}.
\]
Finally, the parameter is updated via
\begin{equation}\label{phibe-q-rule}
\theta_{n+1} = \theta_n +  \frac{\alpha_n}{|\mathcal{B}_n|} \sum_{(s_i,s_i',a_i,r_i)\in\mathcal{B}_n} \Phi(s_i,a_i)\, \delta(s_i,s_i',a_i,r_i),
\end{equation}

\iffalse
the greedy action 
\[
a_{\max}(s_i) = \arg\max_{a'\in \mathcal{A}} \Phi(s_i,a')^\top \theta_n
\]
is computed, and the corresponding maximal Q-value is defined as 
\[
\mathrm{maxQ}(s) = \Phi(s,a_{\max}(s))^\top \theta_n.
\]

Next, \tcr{the drift and diffusion statistics $\hat b(s)$ and $\hat\Sigma(s)$ are replaced by its unbiased estimates:}
\[
\hat b_i(s_i) \approx \frac{s'_i - s_i}{\Delta t}, \quad 
\hat \Sigma_i(s_i) \approx \frac{(s'_i - s_i)(s'_i - s_i)^\top}{\Delta t}.
\]
These estimates are then used to compute the discrete-time PhiBE operator applied to $\mathrm{maxQ}(s)$:
\[
(\widehat{\mathcal{L}}\,\mathrm{maxQ})(s)
= \hat b_i(s) \cdot \nabla_s (\Phi(s,a_{\max}(s))^\top \theta_n)
+ \frac{1}{2} \hat\Sigma_i(s) : \nabla_s^2 (\Phi(s,a_{\max}(s))^\top \theta_n).
\]
Finally, the gradient is averaged over the mini-batch:
\[
g_n = \frac{1}{|\mathcal{B}_n|} \sum_{(s,a,r)\in\mathcal{B}_n} \Phi(s,a)\, \delta(s,a,r),
\]
and the parameter is updated via a standard stochastic gradient step:
\[
\theta_{n+1} = \theta_n + \alpha_n \, g_n.
\]
\fi


This procedure is repeated until a satisfactory approximation of the Galerkin solution for $Q^*(s,a)$ is obtained. 
A rigorous analysis of the sample complexity is left for future work.
The algorithm is presented in Algorithm~\ref{algo2}.

\begin{algorithm}
\caption{PhiBE Q-learning Method}
\begin{algorithmic}\label{algo2}
\STATE \textbf{Input:} Batch size $|\mathcal{B}_n|$, initial value $\theta_0$, discrete time step $\Delta t$, discount $\beta$, state-next-state-reward-action tuples $\mathcal{B} = \{(s_i, s'_i, r_i,a_i)\}_{i=1}^N$, basis $\Phi(s,a)=(\phi_1,\ldots,\phi_n)^\top$, steps $N_{\text{iter}}$, step-sizes $\{\alpha_n\}$
\STATE \textbf{Output:} Estimated parameter $\theta_{N_{\text{iter}}}$ and Q-function $\hat Q(s,a)=\theta^\top\Phi(s,a)$
\STATE $\th \gets \th_0$
\FOR{$n=0,\dots,N_{\text{iter}}-1$}
%\STATE Compute greedy action: $a_{\max}(s) \gets \arg\max_{a'\in\mathcal{A}} \Phi(s,a')^\top\theta$
    \STATE Compute the optimal state-value function $V(s) \gets \max_a \Phi(s,a)^\top \th$
  \STATE Sample a mini-batch $\mathcal{B}_n\subset\mathcal{B}$, and compute PhiBE-Q error for each $(s_i,s_i',a_i,r_i)\in\mathcal{B}_n$
    \[
      \delta(s_i,s_i',a_i,r_i)= \frac{s_i'-s_i}{\dt}\cdot\nabla_sV(s_i)
        + \tfrac{1}{2\dt}(s_i' - s_i)(s_i' - s_i)^\top:\nabla_s^2V(s_i) - \beta \hat{Q}(s_i,a_i)
    \]
  \STATE Update parameter:
   $\theta\gets\theta +\frac{ \alpha_n}{|\mathcal{B}_n|}\sum_{(s_i,s_i',a_i,r_i)\in\mathcal{B}_n}
    \Phi(s_i,a_i)\,\delta(s_i,a_i,r_i)$
\ENDFOR
\STATE \textbf{return} $\hat Q(s,a)=\theta_{N_{\text{iter}}}^\top\Phi(s,a)$
\end{algorithmic}
\end{algorithm}


Several remarks on the Algorithm are in order. 

First, the algorithm can be viewed as a continuous-time analogue of classical Q-learning \citep{watkins1992q}. The classical approach constructs the TD error as the residual of the discrete-time Bellman equation \eqref{discrete-time be}, whereas our method is based on the PhiBE-Q equation \eqref{pde of Qhatstar}.

Second, the algorithm naturally extends to nonlinear function approximation, such as Deep Neural Networks (DNNs). One constructs the same TD error as in \eqref{td-data} using the parametric form  $Q(s,a;\th_n)$, and updates the parameters via
\[
\th_{n+1} = \th_n + \alpha_n \, \delta(s_i,s_i',a_i,r_i)\nabla_\th Q(s_i,a_i;\th_n).
\]
As in classical RL, convergence guarantees are generally unavailable under nonlinear function approximation. In practice, stability can be improved by decoupling the target from the online network. For example, following \citep{mnih2015human}, one may use a target network and define $V_n = \max_a Q(s,a;\th^-)$,
where $\th^-$ denotes parameters from a delayed copy of the Q-network.

Third, with the continuous-time formulations of $Q^*(s,a)$ in \eqref{pde of Qstar} and $Q^\pi(s,a)$ in \eqref{pde of q}, standard off-policy algorithms including Q-learning \citep{watkins1992q}, DQN \citep{mnih2015human}, TD3-BC \citep{fujimoto2021minimalist}, CQL \citep{kumar2020conservative}, etc can be directly adapted to solve continuous-time RL problems.





Finally, we would like to mention that this algorithm is typically applied in settings where $\max_a \Phi(s,a)^\top \theta$ can be computed efficiently. The numerical experiment in Section~\ref{sec:lqr} provides one such example. For discrete action spaces, $\max_a \Phi(s,a)^\top \theta$ is easy to evaluate; however, $V_n(s)$ is not necessarily differentiable. To address this, one may approximate the greedy action $a^*(s) = \arg\max_a Q(s,a)$ using a softmax policy $\pi(a \mid s) \propto \exp\!\big(Q(s,a)/\lambda\big)$, where $\lambda$ controls the degree of exploration. This leads to a softmax variant of PhiBE-Q-learning. Alternatively, one can approximate $\pt_{s_k}V(s_i) \approx \frac1h(V(s_i + he_k) - V(s_i))$, where $e_k$ denotes the $k$th standard basis vector in the state space and $h>0$ is a small finite-difference step size. Provided that $h<\Delta t$, the additional approximation error introduced by this finite-difference step does not dominate the overall discretization error.
%For continuous action space, when $\Phi(s,a)^\top \th$ is a convex function in $a$ for all $s$, it is usually possible to write the explicit expression of $V_n(s)$. For example, when the basis in the action space is spanned by $1, a, a^2$ and the state space is spanned by $\Phi(x)$, then $V_n(s) = ....$. 
%Add another algorithm, and test it on Cartpole experiment.


\section{Experiments} \label{sec:4.3}
\subsection{LQR with stochastic dynamics}\label{sec:lqr}
We consider the linear dynamics with quadratic reward setting,
\[
b(s,a) = As + Ba, \quad\sigma(s,a) = \sigma, \quad r(s,a) = s^\top Q s+a^\top Ra,
\]
The system parameters are set as $A = B = 0.1$, $Q = R = -1$, and the discount factor $\beta = 1$ or $\beta = 2$. In the noise-free case, $\sigma = 0$, whereas in the noisy case, $\sigma = 0.2$. 

To generate the dataset, we simulate $I$ independent trajectories $\{(s_i, s'_i, r_i,a_i)\}_{i=1}^N$.
Here, $s \sim U(-1,1), a \sim U(-1,1)$ and $s'$ is the resulting next state after applying action $a$,
\[
\rho(s' \mid s,a) \sim \mathcal{N}\!\left(e^{A \Delta t}s + \frac{B}{A}\!\left(e^{A \Delta t}-1\right)a,\; \frac{\sigma^2}{2A}\!\left(e^{2A \Delta t}-1\right)\right).
\].

We represent $Q_\th(s,a) = \th^\top \Phi(s,a)$ using quadratic features $\Phi(s,a) = (s^2, a^2, 1)$, and initialize the parameter $\theta_0$ randomly. Note that in LQR, the corresponding greedy value function admits the explicit representation
\[
\max_a Q_\theta(s,a)
=
\theta_m^\top \phi(s),\quad \theta_m
=
\left(
\theta_1-\frac{\theta_2^2}{\theta_3},
\,0,\,
\theta_3
\right)^\top.
\]

%In the data-driven setting, we consider mini-batch sizes \[ B \in \{64,128,256,512,10000\}. \]
%Since the action-value function is parameterized as \[ Q_\theta(s,a)=\theta^\top \Phi(s,a), \quad\theta = (\theta_1, \theta_2, \theta_3)^\top \]


The step size in all plots is set to $\alpha_n = 0.1$. Figure~\ref{fig:mainfig3} shows the evolution of the error $|\theta_n - \theta^*|$ for PhiBE under known discrete-time dynamics. We implement the algorithm in \eqref{eq:newalgoite-phibe}, where $\hat b$, $\hat \Sigma$, and the inner products are computed explicitly. The purpose of this experiment is to demonstrate that PhiBE provides an accurate approximation to the continuous-time optimal control problem and to validate our theoretical guarantees.

Figure~\ref{fig:mainfig4} shows the evolution of the error $|\theta_n - \theta^*|$ under discrete-time data, using Algorithm~\ref{algo2}. This experiment demonstrates that the proposed data-driven algorithm performs comparably to the algorithm with known discrete-time dynamics. 




At each iteration, we sample a mini-batch from the dataset, compute the Optimal-PhiBE-Q residual defined in \eqref{td-data} at $\th_n$
\[
\delta(s,a,r) = r + (\widehat{\mathcal{L}}\,\max Q_{\th_n})(s)
+ (1-\beta)\,\max Q_{\th_n}(s) - Q_{\th_n},
\]
and update $\theta_n$ according to the PhiBE-Q-learning rule \eqref{phibe-q-rule}. 

\begin{figure}[H] % [h!] is a common placement option (here, if possible)
    \centering % Centers the entire group of subfigures
    \begin{subfigure}[b]{0.45\textwidth} % [b] aligns them at the bottom, 0.45\textwidth sets the width
        \centering
        \includegraphics[width=\textwidth]{PhiBEQ_deterministic.png} % image1.png will take up the full width of this subfigure environment
        \caption{Deterministic dynamics} % Caption for the first subfigure
        \label{fig:subfig5} % Label for cross-referencing this specific subfigure
    \end{subfigure}
    \hfill % Adds horizontal space between subfigures, distributing them evenly
    \begin{subfigure}[b]{0.45\textwidth}
        \centering
        \includegraphics[width=\textwidth]{PhiBEQ_stochastic.png}
        \caption{Stochastic dynamics}
        \label{fig:subfig6}
    \end{subfigure}
    \caption{
The evolution of PhiBE Q-learning under known discrete-time dynamics. 
Subfigure (a) shows the deterministic setting ($\sigma=0$), while subfigure (b) shows the stochastic setting ($\sigma=0.2$). $\beta = 1$ in the deterministic setting, $\beta = 2$ in the stochastic setting.
In each subfigure, different curves correspond to the exact known-dynamics update and its corresponding PhiBE approximation using different time steps $\Delta t$.
}
    % You can add more subfigures here if needed

    \label{fig:mainfig3} % Label for cross-referencing the main figure
\end{figure}

\begin{figure}[H] % [h!] is a common placement option (here, if possible)
    \centering % Centers the entire group of subfigures
    \begin{subfigure}[b]{0.45\textwidth} % [b] aligns them at the bottom, 0.45\textwidth sets the width
        \centering
        \includegraphics[width=\textwidth]{Q-learning_deterministic_rev1.png} % image1.png will take up the full width of this subfigure environment
        \caption{Deterministic dynamics} % Caption for the first subfigure
        \label{fig:subfig7} % Label for cross-referencing this specific subfigure
    \end{subfigure}
    \hfill % Adds horizontal space between subfigures, distributing them evenly
    \begin{subfigure}[b]{0.45\textwidth}
        \centering
        \includegraphics[width=\textwidth]{Q-learning_stochastic_rev1.png}
        \caption{Stochastic dynamics}
        \label{fig:subfig8}
    \end{subfigure}
    % You can add more subfigures here if needed
\caption{The evolution of PhiBE Q-learning when only discrete trajectory off-policy data are available with $\dt = 0.1$.
Subfigure (a) shows the deterministic setting ($\sigma=0$), while subfigure (b) shows the stochastic setting ($\sigma=0.2$).
In each subfigure, the batch-size curves correspond to data-driven PhiBE Q-learning with different batch sizes used at each iteration, while the additional baseline curve corresponds to PhiBE Q-learning with known discrete-time dynamics.
The discount parameter is $\beta = 1$ in the deterministic setting and $\beta = 2$ in the stochastic setting.}

    \label{fig:mainfig4} % Label for cross-referencing the main figure
\end{figure}
For the known-dynamics setting, the results in Figure \ref{fig:mainfig3} show that PhiBE Q-learning closely tracks the exact discrete-time Q-iteration across both deterministic and stochastic dynamics. When the time step $\Delta t$ decreases, the PhiBE approximation becomes more accurate. In the stochastic case the long-term convergence trend remains nearly identical to the deterministic counterpart. The algorithm ultimately reaches the same accuracy level, confirming that the PhiBE Q-learning update is stable and robust to moderate diffusion noise in the transition dynamics.

For the off-policy data setting in Figure \ref{fig:mainfig4}, PhiBE Q-learning successfully recovers the Q-function from sample trajectories, and its convergence behavior is strongly influenced by the batch size. Larger batches lead to smoother error decay and significantly faster convergence, while very small batches introduce higher variance and slower stabilization. Comparing the deterministic and stochastic environments, we observe that stochasticity causes oscillations during the transient phase. However, by adopting a decreasing learning rate in later iterations, the method avoids persistent oscillation and converges to a similar final error plateau. These results demonstrate that PhiBE Q-learning remains effective even with off-policy data and stochastic dynamics, achieving reliable convergence across different sampling regimes.

\begin{figure}[H] % [h!] is a common placement option (here, if possible)
    \centering % Centers the entire group of subfigures
    \begin{subfigure}[b]{0.45\textwidth} % [b] aligns them at the bottom, 0.45\textwidth sets the width
        \centering
        \includegraphics[width=\textwidth]{Q_learning_RL_deterministic.png} % image1.png will take up the full width of this subfigure environment
        \caption{Deterministic dynamics} % Caption for the first subfigure
        \label{fig:subfig9} % Label for cross-referencing this specific subfigure
    \end{subfigure}
    \hfill % Adds horizontal space between subfigures, distributing them evenly
    \begin{subfigure}[b]{0.45\textwidth}
        \centering
        \includegraphics[width=\textwidth]{Q_learning_RL_stochastic.png}
        \caption{Stochastic dynamics}
        \label{fig:subfig10}
    \end{subfigure}
    % You can add more subfigures here if needed
\caption{The evolution of PhiBE Q-learning when only discrete trajectory off-policy data are available compared with RL Q-Learning under same batch size $=64$.
Subfigure (a) shows the deterministic setting ($\sigma=0$), while subfigure (b) shows the stochastic setting ($\sigma=0.2$).
In each subfigure, different curves correspond to whether the method is based on PhiBE. We set $\Delta t = 0.1$, $\beta = 2$ in this setting.}
    \label{fig:Q} % Label for cross-referencing the main figure
\end{figure}

To further demonstrate the effectiveness of our method, we also evaluate the standard Q-Learning algorithm using the same data and experimental configuration. As shown in Figure~\ref{fig:Q}, the Q-Learning method iterates converge with a larger error. Therefore, the experimental results clearly confirm that our approach is significantly more robust than the standard Q-learning in the continuous setting.
     
\section{Conclusion}
In this work, we propose a continuous-time Q function, enabling off-policy data usage in continuous-time RL settings. By introducing an time-evolutionary equation and leveraging the Galerkin projection, we develop a stable iterative scheme for solving continuous-time optimal control problems in a model-free way. 

We further establish convergence guarantees for the proposed algorithms under linear function approximation. The resulting stability conditions are independent of $\dt$, allowing our method to learn robustly from discrete trajectory data. Numerical experiments demonstrate the effectiveness of the proposed algorithms and highlight their superior convergence behavior compared to standard Q-learning, which may diverge in similar continuous-time settings.

Several promising directions remain open for future investigation. First, from a theoretical perspective, some of the assumptions used to establish the current guarantees may potentially be relaxed. While this paper focuses on convergence and convergence rate analysis, establishing theoretical sample complexity bounds is an important next step. Second, extending our methodology to large-scale problems via deep neural approximators presents both opportunities and challenges in expressiveness and stability. Third, incorporating exploration strategies and safety constraints would broaden applicability in real-world control systems. 

Overall, this work provides a principled and data-driven foundation for continuous-time reinforcement learning for off-policy data, and we believe the proposed framework will facilitate further developments toward scalable and reliable continuous-time decision-making.


\appendix
\section{Proof of Lemma \ref{lemma:uniform-optimal}}\label{sec:proof of uniform-optimal}

\subsection*{Setting}

Fix a state marginal $\mu(s)$.
The set of all Markov kernel behavior policies is
\begin{equation*}
\Pi := \Bigl\{\pi(\cdot|\cdot) : \pi(a|s) \geq 0,\;
\int_\mathcal{A}\pi(a|s)\,da = 1 \quad \mu(s)\text{-a.e.}\Bigr\},
\end{equation*}
and the uniform policy is $\pi^u(a|s) := 1/|\mathcal{A}|$.
The Gram matrix and coverage coefficient under $\pi$ are
\begin{equation*}
G^\pi := \mathbb{E}_s\!\left[\int_\mathcal{A}\pi(a|s)\,\Phi(s,a)\Phi(s,a)^\top\, da\right],
\end{equation*}
\begin{equation*}
c_0(\pi,\Phi)^2 := \max_{v \in \mathbf{S}^{n-1}}
\frac{\mathbb{E}_s\!\left[\max_{a\in\mathcal{A}}(v^\top\Phi(s,a))^2\right]}
     {v^\top G^\pi v}.
\end{equation*}
Since $\mathcal{A}$ is compact and $\Phi$ is continuous, the maximum over $a$
is attained; since $\mathbf{S}^{n-1}$ is compact and $v\mapsto c_0(\pi,\Phi)^2$
is continuous, the maximum over $v$ is also attained.
Define
\begin{equation*}
\mathcal{F}_B := \{\Phi \text{ continuous} : \|\Phi\|_\infty \leq B\},
\qquad
\bar{H}(\Phi) := \mathbb{E}_s\!\left[\int_\mathcal{A}\Phi(s,a)\Phi(s,a)^\top\, da\right].
\end{equation*}




\begin{proof}
We establish the result in two parts.

\paragraph{Part (I): Upper bound on $c_0(\pi^u, \Phi)^2$ for each fixed $\Phi$.}

For any fixed $\Phi \in \mathcal{F}_B$ and any $v \in \mathbf{S}^{n-1}$,
since $\|\Phi\|_\infty \leq B$:
\begin{equation*}
\mathbb{E}_s\!\left[\max_{a \in \mathcal{A}}(v^\top\Phi(s,a))^2\right] \leq B^2.
\end{equation*}
Under $\pi^u$:
\begin{equation*}
v^\top G^u v
= \frac{1}{|\mathcal{A}|}\,\mathbb{E}_s\!\left[\int_\mathcal{A}(v^\top\Phi(s,a))^2\,da\right]
= \frac{v^\top\bar{H}(\Phi)\,v}{|\mathcal{A}|}
\geq \frac{\lambda_{\min}(\bar{H}(\Phi))}{|\mathcal{A}|}.
\end{equation*}
Dividing and taking the maximum over $v$:
\begin{equation}
c_0(\pi^u,\Phi)^2
\leq \frac{B^2\,|\mathcal{A}|}{\lambda_{\min}(\bar{H}(\Phi))}. \label{eq:part1}
\end{equation}

\paragraph{Part (II): Every $\pi \neq \pi^u$ is strictly suboptimal.}

Fix any $\pi \in \Pi$ with $\pi \neq \pi^u$. Define
\begin{equation*}
A_- := \left\{a \in \mathcal{A} :
\mathbb{E}_s[\pi(a|s)] < \frac{1}{|\mathcal{A}|}\right\}.
\end{equation*}

\textit{Claim: $|A_-| > 0$.}
Since $\pi(\cdot|s)$ and $\pi^u$ are both probability densities on $\mathcal{A}$:
\begin{equation*}
\int_\mathcal{A}\mathbb{E}_s[\pi(a|s)]\,da
= \mathbb{E}_s\!\left[\int_\mathcal{A}\pi(a|s)\,da\right] = 1
= \int_\mathcal{A}\frac{1}{|\mathcal{A}|}\,da.
\end{equation*}
So $\mathbb{E}_s[\pi(\cdot|s)] - 1/|\mathcal{A}|$ integrates to zero over
$\mathcal{A}$ but is not identically zero (since $\pi \neq \pi^u$).
Therefore it must be strictly negative on a set of positive Lebesgue measure,
i.e., $|A_-| > 0$.

\textit{Adversarial basis functions.}
Fix any $v_0 \in \mathbf{S}^{n-1}$ and define
\begin{equation*}
\Phi^*(s,a) := B\,v_0\,\mathbf{1}_{A_-}(a).
\end{equation*}
Note that $\Phi^*$ does not depend on $s$ and $\|\Phi^*\|_\infty \leq B$,
so $\Phi^* \in \mathcal{F}_B$.

\textit{Numerator} (same for all $\pi$):
\begin{equation}
\mathbb{E}_s\!\left[\max_{a \in \mathcal{A}}(v_0^\top\Phi^*(s,a))^2\right]
= B^2. \label{eq:num}
\end{equation}

\textit{Denominator under $\pi$:}
\begin{equation}
v_0^\top G^\pi v_0
= \mathbb{E}_s\!\left[\int_\mathcal{A}(v_0^\top\Phi^*(s,a))^2\,\pi(a|s)\,da\right]
= B^2\int_{A_-}\mathbb{E}_s[\pi(a|s)]\,da. \label{eq:den_pi}
\end{equation}

\textit{Denominator under $\pi^u$:}
\begin{equation}
v_0^\top G^u v_0
= B^2\int_{A_-}\frac{1}{|\mathcal{A}|}\,da
= \frac{B^2\,|A_-|}{|\mathcal{A}|}. \label{eq:den_u}
\end{equation}

\textit{Strict inequality.}
By definition of $A_-$, we have
$\mathbb{E}_s[\pi(a|s)] < 1/|\mathcal{A}|$ for every $a \in A_-$, so:
\begin{equation}
\int_{A_-}\mathbb{E}_s[\pi(a|s)]\,da
< \frac{|A_-|}{|\mathcal{A}|}. \label{eq:strict}
\end{equation}

\textit{Conclusion.}
Since the numerator $B^2$ in \eqref{eq:num} is identical for both policies,
combining \eqref{eq:den_pi}--\eqref{eq:strict}:
\begin{equation*}
c_0(\pi,\Phi^*)^2
= \frac{B^2}{v_0^\top G^\pi v_0}
= \frac{1}{\int_{A_-}\mathbb{E}_s[\pi(a|s)]\,da}
\overset{\eqref{eq:strict}}{>}
\frac{|\mathcal{A}|}{|A_-|}
= \frac{B^2}{v_0^\top G^u v_0}
= c_0(\pi^u,\Phi^*)^2.
\end{equation*}
Therefore
\begin{equation*}
\sup_{\Phi\in\mathcal{F}_B}\, c_0(\pi,\Phi)^2
\geq c_0(\pi,\Phi^*)^2
> c_0(\pi^u,\Phi^*)^2,
\end{equation*}
and since $\pi \neq \pi^u$ was arbitrary, $\pi^u$ is the unique minimizer.
\end{proof}

\subsection*{Corollary}

For any $\Phi \in \mathcal{F}_B$ with $\lambda_{\min}(\bar{H}(\Phi)) \geq \sigma > 0$:
\begin{equation*}
\min_{\pi\in\Pi}\;\sup_{\Phi\in\mathcal{F}_B}\, c_0(\pi,\Phi)
= \sup_{\Phi\in\mathcal{F}_B}\, c_0(\pi^u,\Phi)
\leq \frac{B\sqrt{|\mathcal{A}|}}{\sqrt{\sigma}}.
\end{equation*}
Consequently, the convergence condition $\beta \geq Lc_0$ holds under $\pi^u$
whenever
\begin{equation*}
\beta \;\geq\; \frac{LB\sqrt{|\mathcal{A}|}}{\sqrt{\sigma}},
\qquad \sigma := \lambda_{\min}(\bar{H}(\Phi)),
\end{equation*}
and no other $\pi \in \Pi$ guarantees this bound uniformly over $\mathcal{F}_B$.




\bibliography{PDL.bib} 








\end{document}



\begin{itemize}
    \item \textbf{Case 1: Given Discrete-Time Transition Distribution.} In this scenario, the transition distribution $\rho(s', \Delta t | s)$ for a discrete time step $\Delta t$, driven by the continuous dynamics (2), is provided. Here, $\rho(s', \Delta t | s)$ represents the probability density function of $s_{\Delta t}$ given $s_0 = s$. Under this condition, the discrete-time Bellman equation can be explicitly formulated. The error analysis in Section 3 is conducted primarily under this case.

    \item \textbf{Case 2: Given Trajectory Data.} This scenario involves trajectory data generated by the continuous dynamics (2) and collected at discrete time points $j\Delta t$. The trajectory dataset $\mathcal{S} = \{s^l_0, s^l_{\Delta t}, \dots, s^l_{m\Delta t}\}_{l=1}^I$ comprises $I$ independent trajectories, with the initial state $s^l_0$ of each trajectory sampled from a distribution $\rho_0(s)$.
\end{itemize}

 We observe that conventional discrete-time Bellman equations and associated algorithms exhibit inherent approximation errors and limitations when approximating continuous-time value functions. To address these challenges, this paper introduces the \textbf{PDE-based Bellman Equation (PhiBE)}. Building upon PhiBE, we propose novel methodologies: \textbf{PhiBE-Difference Learning} for policy evaluation and \textbf{PhiBE Q-Learning} for optimal control problems, aiming to provide superior approximations to the true value and Q-functions. Subsequent sections will elaborate on PhiBE, detailing its algorithms and theoretical guarantees under both the known discrete-time transition distribution (Section 3) and data-driven (Section 4) scenarios.


\begin{algorithm}
\caption{Policy Evaluation  under LQR setting}
\begin{algorithmic}
\STATE \textbf{Input:} time step $\Delta t$, number of samples $N$, discount factor $\beta$, learning rate $\alpha$, tolerance \texttt{tol}
\STATE \textbf{Output:} Estimated value function parameter $\theta$ and value $\hat V^\pi(s)=\theta^\top\Phi(s)$
\STATE Generate dataset:
\FOR{$t=1$ \TO $N$}
    \STATE Sample initial state $s_0^{(t)} \sim \mathcal{U}(-1,1)$ and action $a_0^{(t)} \sim \pi(a|s_0^{(t)})$
    \STATE Simulate transition:
    \[
    s_{\Delta t}^{(t)} \gets s_0^{(t)} e^{\Delta t} + a_0^{(t)} (e^{\Delta t} - 1)
    \]
    \STATE Store $(s_t, a_t, sdt_t) \gets (s_0^{(t)}, a_0^{(t)}, s_{\Delta t}^{(t)})$
\ENDFOR
\STATE Construct basis matrix $\Phi \in \mathbb{R}^{N \times 2}$ with rows
\[
\Phi[j,:] = \begin{bmatrix} (s_j)^2 & 1 \end{bmatrix}
\]
\STATE Initialize parameter $\theta$
\REPEAT
    \STATE $\delta \gets 0$
    \FOR{$j=1$ \TO $N$}
        \STATE Compute residual:
        \[
        \delta_0 \gets r \cdot \Phi[j,:] - \beta \, \eta \cdot \Phi[j,:] 
        + \frac{sdt_j - s_j}{\Delta t} \cdot 2 \eta^{(1)} s_j
        \]
        \STATE Accumulate: $\delta \gets \delta + \tfrac{1}{N} \delta_0 \cdot \Phi[j,:]$
    \ENDFOR
    \STATE Update:
    \[
    \theta \gets \theta + \alpha (\Phi \Phi^\top)^{-1} \delta
    \]
\UNTIL{$\|\theta - \theta_{\text{prev}}\| < \texttt{tol}$}
\STATE \textbf{return} $\hat V^\pi(s) = \theta^\top \Phi(s)$
\end{algorithmic}
\end{algorithm}


\begin{algorithm}
\caption{Data-driven Method}
\begin{algorithmic}\label{algo1}
\STATE \textbf{Input:} Discrete time step $\Delta t$, discount coefficient $\beta$, trajectory data $\mathcal{B} = \{(s^l_{j\Delta t}, r^l_{j\Delta t})_{j=0}^m\}_{l=1}^I$ generated by policy $\pi$, basis functions $\Phi(s) = (\phi_1(s), \ldots, \phi_n(s))^\top$, order $i$
\STATE \textbf{Output:} Estimated value function $\hat{V}^\pi$
\STATE Compute matrix $A_i$:
\[
A_i = \sum_{l=1}^I \sum_{j=0}^{m-i} \Phi(s^l_{j\Delta t}) \left[ \beta \Phi(s^l_{j\Delta t}) - \hat{b}_i(s^l_{j\Delta t}) \cdot \nabla \Phi(s^l_{j\Delta t}) - \tfrac{1}{2} \hat{\Sigma}_i(s^l_{j\Delta t}) : \nabla^2 \Phi(s^l_{j\Delta t})^\top \right]
\]
\STATE Where:
\[
\hat{b}_i(s^l_{j\Delta t}) = \sum_{k=1}^i a^i_k (s^l_{(j+k)\Delta t} - s^l_{j\Delta t})
\]
\[
\hat{\Sigma}_i(s^l_{j\Delta t}) = \frac{1}{\Delta t} \sum_{k=1}^i a^i_k (s^l_{(j+k)\Delta t} - s^l_{j\Delta t}) (s^l_{(j+k)\Delta t} - s^l_{j\Delta t})^\top
\]
\STATE Compute vector $b_i$:
\[
b_i = \sum_{l=1}^I \sum_{j=0}^{m-i} r^l_{j\Delta t} \Phi(s^l_{j\Delta t})
\]
\STATE Solve linear system: $\theta = A_i^{-1} b_i$
\STATE \textbf{Return:} $\hat{V}^\pi(s) = \theta^\top \Phi(s)$
\end{algorithmic}
\end{algorithm}

\begin{algorithm}
\caption{LQR problem with PhiBE update}
\begin{algorithmic}
\STATE Compute discretized dynamics: 
\[
\hat{A} \gets \frac{1}{\Delta t} \left( e^{A \Delta t} - 1 \right), \quad \hat{B} \gets \frac{\hat{A} \cdot B}{A}
\]
\STATE Initialize $\eta$
\WHILE{$\| \eta - \eta_{\text{prev}}\| \geq \text{tol}$}
    \STATE $\eta_m \gets \text{maxop\_LQR}(\eta)$
    \STATE Compute update:
    \[
    \Delta \eta \gets \alpha \cdot \left( r + \begin{bmatrix} 2\hat{A} \eta_m^{(1)} \\ 2\hat{B} \eta_m^{(1)} \\ 0 \\ \sigma \eta_m^{(1)} \end{bmatrix} + (1 - \beta) \eta_m - \eta \right)
    \]
    \STATE $\eta \gets \eta + \Delta \eta$
\ENDWHILE
\end{algorithmic}
\end{algorithm}


\begin{algorithm}
\caption{Data-driven Method}
\begin{algorithmic}
\STATE Set step size $\Delta t$, number of samples $N$
\STATE Generate dataset:
\FOR{$t = 1$ \TO $N$}
    \STATE Sample $s_0^{(t)} \sim \mathcal{U}(-1,1)$, $a_0^{(t)} \sim \mathcal{U}(-1,1)$
    \STATE Compute $s_{\Delta t}^{(t)} \gets s_0^{(t)} \cdot e^{\Delta t} + a_0^{(t)} \cdot (e^{\Delta t} - 1)$
    \STATE Store $s_t \gets s_0^{(t)}$, $a_t \gets a_0^{(t)}$, $sdt_t \gets s_{\Delta t}^{(t)}$
\ENDFOR


% \STATE Estimate drift term: $\hat{b} \gets (s_{\Delta t} - s)/\Delta t$
\STATE Construct basis matrix $\Phi$:
\[
\Phi[j,:] = \begin{bmatrix}
s_j^2 & s_j \cdot a_j & a_j^2 & 1
\end{bmatrix}
\]

\STATE Initialize $\eta$
\WHILE{$\| \eta - \eta_{\text{prev}}\| \geq \text{tol}$}
    \STATE $\eta_m \gets \text{maxop\_LQR}(\eta)$
    \FOR{$j = 1$ \TO $N$}
        \STATE $\delta_0 \gets r \cdot \Phi[j,:] + (1 - \beta) \eta_m \cdot \Phi[j,:] + \frac{sdt-s}{dt} \cdot 2 \eta_m^{(1)} s_j - \eta \cdot \Phi[j,:]$
        \STATE $\delta \gets \delta + \delta_0 \cdot \Phi[j,:] / N$
    \ENDFOR
    \STATE $\eta \gets \eta + \alpha \cdot (\Phi\Phi^\top)^{-1} \delta$
\ENDWHILE
\end{algorithmic}
\end{algorithm}

We propose an algorithm for policy evaluation in continuous-time reinforcement learning based on the $i$-th order PhiBE-difference equation. 
Given discrete-time trajectory data generated by a fixed policy $\pi$ and a set of basis functions $\Phi(s)$, 
the algorithm estimates the local drift $\hat{b}_i$ and diffusion $\hat{\Sigma}_i$ using $i$-order PhiBE approximation with coefficients $a_k^i$. 
These estimates are substituted into the PhiBE operator to form the empirical matrix $A_i$, 
while the observed rewards define the vector $b_i$. 
The value function is then obtained by solving the linear system $\theta = A_i^{-1} b_i$, 
yielding the approximation $\hat{V}^\pi(s) = \theta^\top \Phi(s)$. 
This procedure directly leverages discrete trajectory data without explicitly modeling the underlying dynamics, 
and the order parameter $i$ controls the discretization accuracy. 

Let $C_1:=c_b+c_\Sigma$. From the error bounds we have
\[
\delta_b+\delta_\Sigma \le C_1\Delta t.
\]
The perturbed left-hand side is
\begin{align*}
\big(A+\delta_b+\delta_\Sigma\big)^2
&= A^2 + 2A(\delta_b+\delta_\Sigma) + (\delta_b+\delta_\Sigma)^2 \\
&\le A^2 + 2A C_1 \Delta t + C_1^2 \Delta t^2.
\end{align*}
Meanwhile, the right-hand side satisfies
\[
\lambda_{\min}(\widehat\Sigma)\ge \lambda_{\min}(\Sigma)-c_\lambda\Delta t,
\]
so that
\[
2\beta\lambda_{\min}(\widehat\Sigma)\ge 2\beta\lambda_{\min}(\Sigma) - 2\beta c_\lambda\Delta t.
\]

Since by assumption $\gamma=2\beta\lambda_{\min}(\Sigma)-A^2>0$, it suffices to ensure
\[
2A C_1 \Delta t + C_1^2 \Delta t^2 + 2\beta c_\lambda\Delta t < \gamma.
\]
Dropping the quadratic term and requiring
\[
\Delta t \le \frac{\gamma}{4A C_1 + 4\beta c_\lambda}
\]
provides a simple sufficient condition (for smaller $\Delta t$, the quadratic term is negligible; a tighter bound can be obtained by solving the quadratic inequality exactly). 
Hence, one may set
\[
\Delta t_0:=\min\!\left\{1,\ \frac{\gamma}{4A C_1 + 4\beta c_\lambda}\right\},
\]
which ensures the perturbed inequality holds for all $\Delta t\le\Delta t_0$. 
This completes the proof.

\section{PhiBE-Difference Learning}
In this chapter, we develop an online method based on the PhiBE, as introduced in Sections~\ref{sec:3.1.1} and \ref{sec:3.1.2}, to solve policy evaluation (PE) problems given the discrete-time transition dynamics $\rho_{\Delta t}(s'|s,a)$ (\ref{def:PE}). We rigorously establish its convergence and provide an analysis of the convergence rate in Theorem~\ref{thr:conv1}. A model-free, data-driven algorithm is then presented in Section~\ref{sec:3.2}. Finally, numerical experiments are conducted in Section~\ref{sec:3.3} to illustrate the performance of the proposed approach.

\subsection{Algorithm and Theoretical Guarantee}
In this section, we assume we have access to transition distribution $\rho_{\Delta t}(s'|s,a)$. Therefore, we omit the hat here because $\hat{b},\hat{\Sigma}$ are known to us according to ($\ref{eq:phibe}$). we will derive the similar result when only trajectory data is available in section \ref{sec:3.2}.

In this section, we work under the following standing assumption:
\begin{assumption}\label{ass:res1}
We assume that
\[
   \left( \|b\|_{L^\infty} + \tfrac{1}{2} \|\nabla\Sigma\|_{L^\infty}\right)^2 
   < 2\beta \lambda_{\min}(\Sigma).
\]
Intuitively, the condition requires that the magnitude of the drift term $b$ and the spatial variation of the diffusion matrix $\Sigma$ be sufficiently small relative to the discount factor $\beta$ and the diffusion itself. That is, when the drift and the spatial variability of $\Sigma$ decrease, while the noise intensity and $\beta$ increase, the assumption is more easily satisfied.

\end{assumption}
\subsubsection{PDE difference}\label{sec:3.1.1}
The equation (\ref{eq:PE}) could also be rewritten as follows
\begin{equation}\label{eq:hjb}
    r(s) + \mL V^\pi(s) - \beta V^\pi(s) = 0.
\end{equation}

Now we define a new function $V(t,s)$ that satisfying the following PDE,
\begin{equation}
\pt_t V(t,s) = r(s) + \mL V(t,s) - \beta V(t,s)\label{ed:PDEdifference},
\end{equation}
as we will prove in the following Lemma that the solution $V(t,s)$ to the above PDE converges to the target value function $V^\pi(s)$ in (\ref{eq:hjb}) exponentially fast in time.

\begin{lemma}
    Let $E(t) = \ll V(t,s)- V^\pi \rl^2$, under assumption \ref{ass:res1}, there exists $c > 0$, such that $E(t)$ satisfies
    \[
    E(t) \leq e^{-2c t}E(0),
    \]
    where $c = \bigg(\frac{(\|b\|_{L^\infty} + \tfrac{1}{2} \|\nabla\Sigma\|_{L^\infty})^2}{2\lam_{\min}(\Sigma)} - \beta \bigg)$
\end{lemma}
\begin{proof}
    Inserting $r(s) = -\mL V^\pi(s) + \beta V^\pi(s)$ into the above equation, one has
\[
\pt_t (V(t,s)- V^\pi ) = \mL (V(t,s)-V^\pi) - \beta (V(t,s) - V^\pi)
\]
Then with $e(t,s) =V(t,s)- V^\pi$ one has
\begin{align*}
    \frac12\pt_t E(t) &= \la\mL e, e\ra - \beta E(t) \\
                      &= \la(b\nabla + \frac{1}{2}\Sigma:\nabla^2) e, e\ra - \beta E(t)\\
                      &= (b - \frac{1}{2}\nabla\Sigma)\la\nabla e, e\ra-\frac{1}{2}\Sigma\la\nabla e, \nabla e\ra-\beta E(t)\\
                      &\leq (\|b\|_{L^\infty} + \tfrac{1}{2} \|\nabla\Sigma\|_{L^\infty})\la\nabla e, e\ra -\frac{1}{2}\lam_{\min}(\Sigma)\la\nabla e, \nabla e\ra-\beta E(t)\\
                      &\leq \bigg(\frac{(\|b\|_{L^\infty} + \tfrac{1}{2} \|\nabla\Sigma\|_{L^\infty})^2}{2\lam_{\min}(\Sigma)} - \beta \bigg)E(t) =:-c(b,\Sigma,\beta)E(t).  
\end{align*}
The third equality follows from integration by parts. The next inequality follows from the basic inequality. The final inequality follows from the assumption \ref{ass:res1}, so $c >0$.
Then it yields
\[
E(t) \leq e^{-2ct}E(0)
\]
\end{proof}
This means that $V$ will converge to $V^\pi$ exponentially fast.
\subsubsection{PhiBE difference} \label{sec:3.1.2}
Now the goal becomes approximating $V(\infty,s) = V^\pi(s)$. Therefore, we first apply the Galerkin method to handle the spatial dependence in the state variable $s$ by projecting the infinite-dimensional PDE onto a finite-dimensional function space. Then we discretize the resulting system in the time direction, which can be solved numerically.

Applying Galerkin method, we represent $V(s,t) = \th(t)^\top\Phi(s)$ using $n$-dimensional bases $\Phi(s):\S\to\R^n$. 
First the Galerkin solution $(\theta^\pi)^\top \Phi(s)$ that approximates $V^\pi(s)$ satisfies, 
\[
\la r(s) + (\mL-\beta) \Phi(s)^\top\th^\pi,\Phi(s)\ra= 0,
\] 
where we define the inner product as
\[
\la f,g\ra := \int_{\mathbb{R}^d} f(s) g(s) \mu(s) ds,
\]
$\mu$ is a measure over $\mathbb{R}^d$.
Adding time dependence on $\theta(t)$, we derive
\[
\frac{d}{dt}\th(t) = \la r(s) + (\mL-\beta) \Phi(s)^\top\th,\Phi(s)\ra
\] 
Discretize the time evolution, one has
\begin{equation}\label{up:PhiBE}
\th_{n+1} =\th_n+\alpha_n\l[\la r(s) + (\mL-\beta) \Phi(s)^\top\th_n,\Phi(s)\ra\r], 
\end{equation}
this is the update rule for the iterative method when $\rho_{\Delta t}(s'|s,a)$ is available.

We next establish a sufficient condition under which the associated iteration converges to Galerkin solution of equation (\ref{eq:hjb}).
\begin{theorem}\label{thr:conv1}
Let $\Phi(s) \in \mathbb{R}^n$ be a smooth basis functions and $\mu$ be a probability measure over $\mathcal{S}$. Suppose assumption \ref{ass:res1} hold, then the matrix
\[
 A_\Phi:= \int_{\mathcal{S} } \Phi \left( b \cdot \nabla \Phi + \tfrac{1}{2} \Sigma:\nabla^2 \Phi - \beta\Phi \right)^\top d\mu
\]
satisfies $\mathrm{Re}(\lambda_i(A_\Phi)) < 0$ for all $i$. Consequently, the iterative method
\[
\theta_{k+1} = \theta_k + \eta (A_\Phi \theta_k - b_\Phi), \quad \text{where } b_\Phi = \int r^\top \Phi \, d\mu.
\]
converges to Galerkin solution of equation (\ref{eq:hjb}) for all $\eta$ s.t. \[
0 < \eta < \min_{\lambda \in \sigma(A_\Phi)} \frac{-2\,\Re(\lambda)}{|\lambda|^2},
\]
With convergence rate 
\[
\|\theta_k - \theta^\ast\| \;\le\; \bigg(\max_{\lambda \in \sigma(A_\Phi)} |1+\eta \lambda|\bigg)^k \,\|\theta_0 - \theta^\ast\|.
\]
As for the relation between arguments and convergence rate, when the drift term $b$ and the spatial variation of $\Sigma$ decrease, while the noise intensity and the discount factor $\beta$ increase, the convergence rate becomes faster.

The restriction is formally stated in terms of the drift and diffusion in the PhiBE setting. However, since the error bounds $\|b - \hat{b}\|_{\infty}$, $\|\Sigma - \hat{\Sigma}\|_{\infty}$, and $\|\nabla\Sigma - \nabla\hat{\Sigma}\|_{\infty}$ vanish as $\Delta t \to 0$, the same restriction can effectively be interpreted as a condition on the original dynamics parameters.

Lastly, in the PhiBE setting, given $\rho_{\Delta t}(s' \mid s,a)$, we have derived the convergence rate of the iterative method. However, it is unrealistic to assume access to infinite data. The analysis of the finite sample case is left for future work.
\end{theorem}


\begin{proof}
We write the linear system as
\[
A_\Phi \theta = b_\Phi. 
\]

To prove convergence of the iteration $\theta_{k+1} = \theta_k + \eta(A_\Phi \theta_k - b_\Phi)$, it suffices to show that all eigenvalues of $A_\Phi$ have negative real parts.

Applying integral by parts, we have
\[
A_\Phi =  \int \Phi \left( b \cdot \nabla \Phi - \tfrac{1}{2} \nabla \Sigma \nabla \Phi  - \beta\Phi \right)^\top - \tfrac{1}{2} \nabla\Phi\Sigma\nabla\Phi^\top d\mu.
\]

We estimate the operator norm of the first two term in $A_\Phi$ using Hölder and Cauchy–Schwarz inequalities:
\[
\left\| \int \Phi \left( b \cdot \nabla \Phi - \tfrac{1}{2} \nabla \Sigma \nabla \Phi \right)^\top d\mu \right\| 
\leq  \|\Phi\|_{L^2(\mu)} \|\nabla \Phi\|_{L^2(\mu)}\left( \|b\|_{L^\infty} + \tfrac{1}{2} \|\nabla\Sigma\|_{L^\infty}\right).
\]
Then 
\begin{align*}
    v^\top (\frac{A_\Phi + A_\Phi^\top}{2}) v &\leq \|v^\top\Phi\|_{L^2(\mu)} \|v^\top\nabla \Phi\|_{L^2(\mu)}\left( \|b\|_{L^\infty} + \tfrac{1}{2} \|\nabla\Sigma\|_{L^\infty}\right) - \beta\|v^\top\Phi\|_{L^2(\mu)}^2-\frac{\lam_{\min}(\Sigma)}{2}\|v^\top\nabla \Phi\|_{L^2(\mu)}^2\\
    &\leq \|v^\top\Phi\|_{L^2(\mu)} \|v^\top\nabla \Phi\|_{L^2(\mu)}\bigg( \|b\|_{L^\infty} + \tfrac{1}{2} \|\nabla\Sigma\|_{L^\infty}-2\sqrt{\frac{\beta\lam_{\min} (\Sigma)}{2}}\bigg).
\end{align*}
With the given assumption, we conclude that:
\[
v^\top (\frac{A_\Phi + A_\Phi^\top}{2}) v < 0 \quad \text{for all } v \ne 0,
\]
so all eigenvalues of $A_\Phi$ lie in the left half-plane. If the spectral radius $\rho(I + \eta A_{\Phi}) < 1$, then the convergence is guaranteed. That is,
\[
|1 + \eta \lambda| < 1, \quad \forall \lambda \in \sigma(A_\Phi),
\]
where $\sigma(A_\Phi)$ denotes all eigenvalues of $A_\Phi$
Since $\Re(\lambda) < 0$ for all eigenvalues, there always exists $\eta > 0$ small enough such that the condition holds. More explicitly, for each eigenvalue $\lambda = -a + i b$ with $a > 0$, the condition reads
\[
(1 - \eta a)^2 + (\eta b)^2 < 1.
\]
This inequality is satisfied whenever $0 < \eta < \frac{2a}{a^2 + b^2}$, and thus the iteration converges for all
\[
0 < \eta < \min_{\lambda \in \sigma(A_\Phi)} \frac{-2\,\Re(\lambda)}{|\lambda|^2}
\]The convergence of the iteration with a suitable step size then follows from standard linear system theory.
\end{proof}

The analysis is based on the known dynamics setting, but with stronger assumption, we could generalize our result to PhiBE setting by following corollary. 
\begin{corollary}\label{cor:PhiBE_short}
Under a strengthened version of Assumption~\ref{ass:res1}
\[
 \left( \|b\|_{L^\infty} + \tfrac{1}{2} \|\nabla\Sigma\|_{L^\infty}\right)^2 
   < 2\beta \lambda_{\min}(\Sigma) - c(b,\Sigma)\Delta t,
\]
where $c(b,\Sigma) = f(c_b,c_{\Sigma}), \|b-\hat b\| \leq c_b\Delta t, \|\Sigma-\hat \Sigma\| \leq c_\Sigma\Delta t$
the conclusion of the theorem proved for the known-dynamics setting remains valid in the PhiBE case.
\end{corollary}




\subsection{Algorithm(Data-driven)}\label{sec:3.2}



In the previous section, we assumed that the discrete-time transition dynamics $\rho_{\Delta t}(s'|s,a)$ is known. Here, we describe how to implement the update rule \eqref{up:PhiBE} in a data-driven manner when the transition dynamics is unknown and only the state–next-state–reward tuples $\{(s_i, s'_i,r_i)\}_{i=1}^n$ are available, where $s_i \sim \mu(s), s'_i \sim \rho_{\Delta t}(s'|s,\pi(s)), r_i = r(s_i,\pi(s_i))$.

Since
\[
\E_{s_i \sim \mu} \Big[\Phi(s_i) \big[(\hat{b}(s_i)\nabla + \hat{\Sigma}(s_i)\nabla^2 - \beta) \Phi(s_i) + r(s_i) \big] \Big] 
= \la r(s) + (\mathcal{L}_{\hat b, \hat\Sigma}-\beta) \Phi(s)^\top \theta_n, \Phi(s)\ra,
\]
it follows that
\[
\Phi(s_i) \Big[ (\hat{b}(s_i)\nabla + \hat{\Sigma}(s_i)\nabla^2 - \beta)\Phi(s_i) + r(s_i) \Big]
\]
is an unbiased estimator of $\la r(s) + (\mathcal{L}_{\hat{b},\hat{\Sigma}}-\beta) \Phi(s)^\top \theta_n, \Phi(s)\ra$. Moreover, by the fact that $s'_i \sim \rho_{\Delta t}(s'|s,\pi(s))$ and the definition of $\hat b(s_i), \hat\Sigma(s_i)$, we know that
\[
\E\bigg [\frac{s'_i - s_i}{\Delta t}\bigg|s_i\bigg] = \hat{b}(s_i) \quad \text{and} \quad \E\bigg [\frac{(s'_i - s_i)(s'_i - s_i)^\top}{\Delta t}\bigg|s_i\bigg] = \hat\Sigma(s_i),
\]
so $\hat{b}_i = \frac{s'_i - s_i}{\Delta t}, \hat \Sigma_i = \frac{(s'_i - s_i)(s'_i - s_i)^\top}{\Delta t}$
are unbiased estimators of $\hat b(s_i)$ and $\hat \Sigma(s_i)$, respectively. Therefore the TD-error provides an unbiased estimate of $\la r(s) + (\mL-\beta) \Phi(s)^\top \theta_n, \Phi(s)\ra$.

We define the update rule (\ref{it:PhiBE}) as a model-free online algorithm solving for the PhiBE \ref{eq:hjb} under linear approximation. In each iteration, the coefficient vector $\theta_k$ associated with the basis $\Phi(s)$ is updated until a satisfactory approximation of the Galerkin solution of (\ref{up:PhiBE}) is obtained. The reason we call it model-free is that its update rule does not depend on the specific dynamics parameters. In other words, regardless of how the underlying dynamics vary, the iterative scheme used to compute the updates remains unchanged. A rigorous analysis of the sample complexity is left for future work.

The algorithm is presented in Algorithm~\ref{algo1}.(\citep{zhu2025optimalphibepdebasedmodelfreeframework}) 


\begin{algorithm}
\caption{Data-driven Iterative Method (one-step PhiBE)}
\begin{algorithmic}\label{algo1}
\STATE \textbf{Input:} Batch size $|\mathcal{B}_n|$, initial value $\theta_0$, discrete time step $\Delta t$, discount $\beta$, 
state-next-state-reward tuples $\mathcal{B} = \{(s_i, s'_i, r_i)\}_{i=1}^N$, basis $\Phi(s) = (\phi_1,\ldots,\phi_n)^\top$, steps $N_{\text{iter}}$, step-sizes $\{\alpha_n\}$
\STATE \textbf{Output:} Estimated parameter $\theta_{N_{\text{iter}}}$ and value $\hat V^\pi(s) = \theta^\top \Phi(s)$
\STATE Initialize $\theta_0$ (e.g. $0$)
\FOR{$n=0,\dots,N_{\text{iter}}-1$}
  \STATE Sample a mini-batch $\mathcal{B}_n \subset \mathcal{B}$
  \STATE Compute one-step PhiBE differences for each pair:
  \[
    \hat b_i(s_i) = \frac{s'_i - s_i}{\Delta t}, \quad
    \hat \Sigma_i(s_i) = \frac{(s'_i - s_i)(s'_i - s_i)^\top}{\Delta t}
  \]
  \STATE Compute empirical operator action:
  \[
    (\widehat{\mathcal{L}}\Phi)(s_i) = \hat b_i(s_i) \cdot \nabla \Phi(s_i) + \tfrac{1}{2} \hat \Sigma_i(s_i) : \nabla^2 \Phi(s_i)
  \]
  \STATE Compute TD-error:
  \[
    g_n = \frac{1}{|\mathcal{B}_n|} \sum_{(s_i,s'_i,r_i)\in \mathcal{B}_n} \Phi(s_i) \Big[r_i + (\widehat{\mathcal{L}} - \beta)\Phi(s_i)^\top \theta_n \Big]
  \]
  \STATE Update parameter:
  \begin{equation}\label{it:PhiBE}
      \theta_{n+1} = \theta_n + \alpha_n g_n
  \end{equation}
\ENDFOR
\STATE \textbf{return} $\hat V^\pi(s) = \theta_{N_{\text{iter}}}^\top \Phi(s)$
\end{algorithmic}
\end{algorithm}


\subsubsection{Trigonometric function basis} 
In this setting, we consider the basis $\sin nx$ and $\cos nx$. $V(s)$ is fixed as $\cos^3(s)$, the reward function is derived as  $\beta \cos(ks)^3 - \lam s(-3k\cos^3(ks)\sin(ks))$.

The value function parameters converge to the analytical solution $V^\pi$ within a small number of iterations, as shown in the error plot.
\begin{figure}[H] % [h!] is a common placement option (here, if possible)
    \centering % Centers the entire group of subfigures
    \begin{subfigure}[b]{0.45\textwidth} % [b] aligns them at the bottom, 0.45\textwidth sets the width
        \centering
        \includegraphics[width=\textwidth]{data_driven_deterministic.png} % image1.png will take up the full width of this subfigure environment
        \caption{Data-driven} % Caption for the first subfigure
        \label{fig:subfig3} % Label for cross-referencing this specific subfigure
    \end{subfigure}
    \hfill % Adds horizontal space between subfigures, distributing them evenly
    \begin{subfigure}[b]{0.45\textwidth}
        \centering
        \includegraphics[width=\textwidth]{PhiBE Known Dynamics.jpg}
        \caption{Known dynamics}
        \label{fig:subfig4}
    \end{subfigure}
    % You can add more subfigures here if needed
    \caption{1st-order PhiBE and 2nd-order PhiBE:$\beta = 10, b(s) = 0.05s$, 10000 iterations.} % Main caption for the entire figure
    \label{fig:mainfig2} % Label for cross-referencing the main figure
\end{figure}
\subsubsection{PhiBE Difference learning vs discrete Bellman equation}
